% (User input window)
%
% 本段将 logistic 尺度核的平移族组织为一组可直接嵌入层析与判别论的闭式字典：
% 首先给出若干核心 \(f\)-散度的完全显式表达及其局部常数链；
% 随后把该核诱导的多类别 MAP 判决、最优布局、互信息与 Fisher 亏损写成可审计公式。

在命题 \ref{prop:xi-bulk-curvature-logistic-kernel-decomposition} 的尺度曲率接口中，深度原子以 logistic 核的平移族出现。
因此，平移族 \(\{\phi_a\}_{a\in\RR}\) 的散度、判决误差与信息几何常数可被视为“尺度可见性阈值”的原型账本。

\begin{proposition}[尺度曲率核的散度字典：\texorpdfstring{$1/4$}{1/4}、\texorpdfstring{$4$}{4}、\texorpdfstring{$1/3$}{1/3} 的同源封口]\label{prop:xi-logistic-divergence-dictionary}
\leanverified{paper\_xi\_logistic\_divergence\_dictionary}
令
\[
\phi(s)=\frac{1}{4\cosh^{2}(s/2)}=\frac{e^{s}}{(1+e^{s})^{2}},
\qquad
\phi_a(s):=\phi(s-a),
\qquad
\Delta:=b-a.
\]
则对任意 \(a,b\in\RR\) 下列散度均存在完全闭式：
\begin{enumerate}
  \item \textup{（总变差距离）}
  \[
  \|\phi_a-\phi_b\|_{\mathrm{TV}}
  :=\frac12\int_{\RR}|\phi_a(s)-\phi_b(s)|\,\dd s
  =\tanh\!\Bigl(\frac{|\Delta|}{4}\Bigr).
  \]
  \item \textup{（二元等先验 Bayes 最小错误率）}
  \[
  P_{\mathrm e}(\Delta)
  =\frac{1-\|\phi_a-\phi_b\|_{\mathrm{TV}}}{2}
  =\frac{1}{1+e^{|\Delta|/2}}.
  \]
  \item \textup{（Hellinger 亲和度与平方 Hellinger 距离）}
  \[
  \int_{\RR}\sqrt{\phi_a(s)\phi_b(s)}\,\dd s
  =\frac{|\Delta|}{2\sinh(|\Delta|/2)},
  \qquad
  H^{2}(\phi_a,\phi_b)
  :=1-\int_{\RR}\sqrt{\phi_a\phi_b}
  =1-\frac{|\Delta|}{2\sinh(|\Delta|/2)}.
  \]
  \item \textup{（Kullback--Leibler 散度的严格对称性）}
  \[
  D_{\mathrm{KL}}(\phi_a\|\phi_b)
  :=\int_{\RR}\phi_a(s)\log\frac{\phi_a(s)}{\phi_b(s)}\,\dd s
  =|\Delta|\coth(|\Delta|/2)-2,
  \qquad
  D_{\mathrm{KL}}(\phi_a\|\phi_b)=D_{\mathrm{KL}}(\phi_b\|\phi_a).
  \]
  \item \textup{（\(\chi^2\) 散度与 Rényi-\(2\) 散度）}
  \[
  \chi^2(\phi_a\|\phi_b)
  :=\int_{\RR}\frac{(\phi_a-\phi_b)^2}{\phi_b}\,\dd s
  =\frac{4}{3}\sinh^{2}(|\Delta|/2),
  \qquad
  D_2(\phi_a\|\phi_b)
  :=\log\!\int_{\RR}\frac{\phi_a(s)^2}{\phi_b(s)}\,\dd s
  =\log\!\Bigl(\frac{2\cosh|\Delta|+1}{3}\Bigr).
  \]
\end{enumerate}
并且在 \(|\Delta|\downarrow 0\) 时存在统一的局部量子化展开
\[
\|\phi_a-\phi_b\|_{\mathrm{TV}}=\frac{|\Delta|}{4}+O(|\Delta|^{3}),
\qquad
H^{2}(\phi_a,\phi_b)=\frac{\Delta^{2}}{24}+O(\Delta^{4}),
\qquad
D_{\mathrm{KL}}(\phi_a\|\phi_b)=\frac{\Delta^{2}}{6}+O(\Delta^{4}),
\]
从而出现严格的局部比例律
\[
D_{\mathrm{KL}}(\phi_a\|\phi_b)=4\,H^{2}(\phi_a,\phi_b)+O(\Delta^{4}),
\qquad
\|\phi_a-\phi_b\|_{\mathrm{TV}}\sim\frac{|\Delta|}{4}.
\]
该常数链把后续尺度可见性阈值固定到同一组不可消去的普适常数上。
\end{proposition}
\begin{proof}
由平移不变性，只需取 \(a=0\)、\(b=\Delta>0\)。
记 \(F(s)=\int_{-\infty}^{s}\phi(u)\dd u=(1+e^{-s})^{-1}\) 为 logistic 分布函数。
因为 \(\phi\) 为偶函数且关于 \(0\) 单峰，\(\phi(s)\) 与 \(\phi(s-\Delta)\) 的唯一交点为 \(s=\Delta/2\)，故
\[
\|\phi-\phi_{\Delta}\|_{\mathrm{TV}}
=\int_{\Delta/2}^{\infty}\bigl(\phi(s-\Delta)-\phi(s)\bigr)\dd s
=F(\Delta/2)-F(-\Delta/2)
=\tanh(\Delta/4),
\]
并由二元 Bayes 公式得到 \(P_{\mathrm e}(\Delta)\)。
\par
对 Hellinger 亲和度，写 \(\phi(s)=e^{s}(1+e^{s})^{-2}\) 并作代换 \(x=e^{s}\)（\(\dd s=\dd x/x\)）得
\[
\int_{\RR}\sqrt{\phi(s)\phi(s-\Delta)}\,\dd s
=e^{-\Delta/2}\int_{0}^{\infty}\frac{1}{(1+x)(1+e^{-\Delta}x)}\,\dd x
=\frac{\Delta}{2\sinh(\Delta/2)}.
\]
对 \(D_{\mathrm{KL}}\)，同样代换得到
\[
D_{\mathrm{KL}}(\phi\|\phi_{\Delta})
=\int_{0}^{\infty}\frac{\Delta+2\log(1+e^{-\Delta}x)-2\log(1+x)}{(1+x)^{2}}\,\dd x.
\]
对
\(
I(\alpha):=\int_{0}^{\infty}\frac{\log(1+\alpha x)}{(1+x)^{2}}\,\dd x
\)
作分部积分可得
\(
I(\alpha)=\alpha\int_{0}^{\infty}\frac{1}{(1+x)(1+\alpha x)}\dd x
=\frac{\alpha\log(1/\alpha)}{1-\alpha}
\)；
取 \(\alpha=e^{-\Delta}\) 并用 \(I(1)=1\) 得
\[
D_{\mathrm{KL}}(\phi\|\phi_{\Delta})
=\Delta+2(I(e^{-\Delta})-I(1))
=\Delta\coth(\Delta/2)-2,
\]
从而对称性由 \(\Delta\mapsto-\Delta\) 的偶性得到。
\par
对 \(\chi^2\) 与 \(D_2\)，由
\(
\chi^2(\phi\|\phi_{\Delta})=\int \phi^2/\phi_{\Delta}-1
\)，同样代换得到
\[
\int_{\RR}\frac{\phi(s)^2}{\phi(s-\Delta)}\,\dd s
=e^{\Delta}\int_{0}^{\infty}\frac{(1+e^{-\Delta}x)^{2}}{(1+x)^{4}}\,\dd x
=\frac{e^{\Delta}+1+e^{-\Delta}}{3}
=\frac{2\cosh\Delta+1}{3},
\]
其中最后一步使用 Beta 积分
\(
\int_{0}^{\infty}\frac{x^{m}}{(1+x)^{4}}\dd x=B(m+1,3-m)
\)
（\(m=0,1,2\)）。
因此
\(
\chi^2(\phi\|\phi_{\Delta})=\frac{2}{3}(\cosh\Delta-1)=\frac{4}{3}\sinh^{2}(\Delta/2)
\)
且 \(D_2=\log(1+\chi^2)\)。
局部展开由 \(\tanh\)、\(\sinh\)、\(\coth\) 的 Taylor 展开直接得到。证毕。
\end{proof}

% (User input window)

\begin{proposition}[Hellinger--Gram 体积势：平方根嵌入的闭式核与碰撞发散律]\label{prop:xi-logistic-hellinger-gram-determinant}
沿用命题 \ref{prop:xi-logistic-divergence-dictionary} 的记号并取深度谱 \(\sigma_1,\dots,\sigma_\kappa\in\RR\) 两两不同。
令
\[
\psi_j(s):=\sqrt{\phi(s-\sigma_j)}\in L^2(\RR),
\qquad
G^{(1/2)}_{jk}:=\langle \psi_j,\psi_k\rangle=\int_{\RR}\sqrt{\phi(s-\sigma_j)\phi(s-\sigma_k)}\,\dd s.
\]
则 \(G^{(1/2)}\) 正定，并且其核具有完全闭式
\[
\boxed{\;
G^{(1/2)}_{jk}=\frac{|\sigma_j-\sigma_k|}{2\sinh\!\bigl(|\sigma_j-\sigma_k|/2\bigr)}\;}
\qquad(1\le j,k\le\kappa),
\]
其中对角线取连续延拓值 \(G^{(1/2)}_{jj}=1\)。
定义体积与熵势
\[
V_{1/2}:=\sqrt{\det G^{(1/2)}},
\qquad
S_{1/2}:=-\log\det G^{(1/2)}.
\]
当发生尺度碰撞 \(\sigma_j=\sigma_\star+\varepsilon t_j\) 且 \(\varepsilon\downarrow 0\)（\(t_j\) 两两不同）时，有普适 Vandermonde 发散律
\[
\boxed{\;
S_{1/2}(\boldsymbol{\sigma})
=\kappa(\kappa-1)\log\frac{1}{\varepsilon}
+2\sum_{1\le j<k\le \kappa}\log\frac{1}{|t_k-t_j|}
-\log\Bigl(\mathfrak H_\kappa\prod_{j=0}^{\kappa-1}\frac{1}{(j!)^{2}}\Bigr)
+o(1).\;}
\]
其中 \(\mathfrak H_\kappa\) 为定理 \ref{thm:xi-hellinger-gram-collision-hankel-constant} 所定义的普适 Hankel 常数（仅依赖于 \(\kappa\) 与谱密度 \(w(\xi)=\pi^2\mathrm{sech}^2(\pi\xi)\)）。
其中二体对数项系数恒为 \(2\)，与谱点权重及其余构型细节无关。
\end{proposition}
\begin{proof}
正定性由 \(G^{(1/2)}\) 的 Gram 表示直接得到。
闭式由命题 \ref{prop:xi-logistic-divergence-dictionary} 第 \textup{(3)} 条（Hellinger 亲和度闭式）取 \(a=\sigma_j\)、\(b=\sigma_k\) 立即推出。
\par
对碰撞渐近，定理 \ref{thm:xi-hellinger-gram-collision-hankel-constant} 给出更精确的行列式渐近：
\[
\det G^{(1/2)}
=\varepsilon^{\kappa(\kappa-1)}
\Bigl(\prod_{j=0}^{\kappa-1}\frac{1}{(j!)^{2}}\Bigr)\,
\mathfrak H_\kappa\,
\prod_{1\le j<k\le\kappa}(t_k-t_j)^{2}\,\bigl(1+o(1)\bigr).
\]
取负对数即得所示展开。证毕。
\end{proof}

% (User input window)
%
% 本段将 Hellinger 二体核 \(A(\Delta)\) 由闭式亲和度提升为一条完全谱确定的几何对象：
% 其 Fourier 密度为 \(\cosh^{-2}(\pi\xi)\)，并诱导一个显式 RKHS 与行列式体积势；
% 在等距谱下进一步 Toeplitz 化并给出 Szeg\H{o} 极限与 CUE 表示。

\begin{proposition}[Hellinger 平移核的 Fourier 完全闭式：\texorpdfstring{$\cosh^{-2}$}{cosh^{-2}} 为深度二体几何的谱密度]\label{prop:xi-hellinger-kernel-fourier-sech2}
\leanverified{paper\_xi\_hellinger\_kernel\_fourier\_sech2}
令尺度单峰核
\[
\phi(s)=\frac{1}{4\cosh^{2}(s/2)}\qquad(s\in\RR),
\]
并定义 Hellinger 亲和核
\[
A(\Delta):=\int_{\RR}\sqrt{\phi(s)\phi(s-\Delta)}\,\dd s
=\frac{|\Delta|}{2\sinh(|\Delta|/2)}\qquad(\Delta\in\RR).
\]
则 \(A\in L^{1}(\RR)\cap C^\infty(\RR)\) 为偶函数。
在 Fourier 变换约定
\(
\widehat f(\xi):=\int_{\RR}e^{-\mathrm{i}\xi x}f(x)\,\dd x
\)
下，有严格恒等式
\[
\boxed{\;
\widehat A(\xi)=\int_{\RR}A(\Delta)e^{-\mathrm{i}\xi\Delta}\,\dd\Delta
=\frac{\pi^{2}}{\cosh^{2}(\pi\xi)}\;}
\qquad(\xi\in\RR).
\]
特别地
\(
\int_{\RR}A(\Delta)\,\dd\Delta=\widehat A(0)=\pi^{2}
\)。
\end{proposition}
\begin{proof}
令
\(
f(s):=\sqrt{\phi(s)}=\frac{1}{2\cosh(s/2)}
\)。
则
\(
A(\Delta)=\int_{\RR}f(s)f(s-\Delta)\,\dd s
\)
为 \(f\) 的自相关，因而 \(\widehat A(\xi)=|\widehat f(\xi)|^{2}\)。
作代换 \(s=2u\) 得
\[
\widehat f(\xi)
=\int_{\RR}\frac{1}{2\cosh(s/2)}e^{-\mathrm{i}\xi s}\,\dd s
=\int_{\RR}\frac{1}{\cosh u}\,e^{-2\mathrm{i}\xi u}\,\dd u
=\frac{\pi}{\cosh(\pi\xi)},
\]
其中最后一步为 \(\cosh^{-1}\) 的 Fourier 闭式。
故 \(\widehat A(\xi)=\pi^{2}/\cosh^{2}(\pi\xi)\)。证毕。
\end{proof}

% (User input window)
%
% 本文件给出 Hellinger 核的 Green 算子刻画与点源能量/电势对偶。

\begin{proposition}[Hellinger 核的 Green 算子刻画：\texorpdfstring{$\cosh^{2}(\pi\sqrt{-\partial^{2}})$}{cosh^2(pi sqrt(-d^2))} 的无限阶椭圆性]\label{prop:xi-hellinger-green-operator-cosh2}
\leanverified{paper\_xi\_hellinger\_green\_operator\_cosh2}
令连续卷积核
\(
A(\Delta)=\frac{|\Delta|}{2\sinh(|\Delta|/2)}
\)
并沿用命题 \ref{prop:xi-hellinger-kernel-fourier-sech2} 的 Fourier 乘子
\[
\widehat A(\xi)=\frac{\pi^2}{\cosh^2(\pi\xi)}.
\]
定义无限阶算子
\[
\mathcal L:=\frac{1}{\pi^2}\cosh^{2}\!\bigl(\pi\sqrt{-\partial_s^{2}}\bigr)
\quad\Longleftrightarrow\quad
\widehat{(\mathcal L f)}(\xi)=\frac{\cosh^{2}(\pi\xi)}{\pi^2}\,\widehat f(\xi).
\]
则在温和分布意义下严格有
\[
\boxed{\;\mathcal L(A)=\delta_0\;},
\]
亦即 \(A\) 为 \(\mathcal L^{-1}\) 的 Green 核。
并且 \(\cosh^{2}(\pi\xi)\sim \tfrac14 e^{2\pi|\xi|}\)（\(|\xi|\to\infty\)）给出高频端的指数放大律，是后续可逆性预算指数爆炸的解析来源。
\end{proposition}
\begin{proof}
由 Fourier 乘子定义，
\[
\widehat{(\mathcal L A)}(\xi)
=\frac{\cosh^{2}(\pi\xi)}{\pi^2}\,\widehat A(\xi)
=\frac{\cosh^{2}(\pi\xi)}{\pi^2}\cdot\frac{\pi^2}{\cosh^{2}(\pi\xi)}
=1.
\]
由于 \(\widehat{\delta_0}\equiv 1\)，故 \(\mathcal L(A)=\delta_0\)（\(\mathcal S'(\RR)\) 中）。证毕。
\end{proof}

\begin{proposition}[点荷—电势对偶：\texorpdfstring{$G$}{G} 与 \texorpdfstring{$G^{-1}$}{G^{-1}} 分别是能量与 Dirichlet-to-Neumann 矩阵]\label{prop:xi-hellinger-charge-potential-duality-dtn}
取两两不同的深度点 \(\{s_j\}_{j=1}^\kappa\subset\RR\)，并令 \(G\in\RR^{\kappa\times\kappa}\) 为 Hellinger--Gram 矩阵
\[
G_{ij}=A(s_i-s_j).
\]
对任意荷载 \(c\in\RR^\kappa\)，定义势函数
\[
u_c(s):=\sum_{j=1}^{\kappa}c_j\,A(s-s_j).
\]
则在命题 \ref{prop:xi-hellinger-green-operator-cosh2} 的 Green 结构下，
\[
\mathcal L u_c=\sum_{j=1}^{\kappa}c_j\,\delta_{s_j}
\]
（分布意义），并且能量严格满足
\[
\boxed{\;
\langle u_c,\mathcal L u_c\rangle
=\sum_{i=1}^{\kappa}c_i\,u_c(s_i)
=c^{\mathsf T}G\,c\;}
\]
其中 \(\langle\cdot,\cdot\rangle\) 为 \(\mathcal S'(\RR)\times\mathcal S(\RR)\) 的自然配对。
反之，若给定结点电势 \(v\in\RR^\kappa\)（\(v_i=u(s_i)\)），则唯一荷载为 \(c=G^{-1}v\)，且对应能量为
\[
\boxed{\;v^{\mathsf T}G^{-1}v\;}
\]
从而 \(G\) 与 \(G^{-1}\) 分别刻画“点源能量”与“点集 Dirichlet-to-Neumann”映射。
\end{proposition}
\begin{proof}
由命题 \ref{prop:xi-hellinger-green-operator-cosh2} 的平移不变性，\(\mathcal L(A(\cdot-s_j))=\delta_{s_j}\)，线性叠加得到 \(\mathcal L u_c=\sum_j c_j\delta_{s_j}\)。
故
\[
\langle u_c,\mathcal L u_c\rangle
=\sum_{i=1}^{\kappa}c_i\,u_c(s_i)
=\sum_{i,j=1}^{\kappa}c_i c_j\,A(s_i-s_j)
=c^{\mathsf T}G\,c.
\]
若 \(v_i=u(s_i)\)，则 \(v=Gc\)；因 \(G\succ0\)（点集两两不同且核严格正定），解唯一且 \(c=G^{-1}v\)。
代回能量等式得到 \(c^{\mathsf T}G c=v^{\mathsf T}G^{-1}v\)。证毕。
\leanverified{paper\_xi\_hellinger\_charge\_potential\_duality\_dtn}
\end{proof}


\begin{proposition}[深度谱的 RKHS 化：\texorpdfstring{$A$}{A} 为严格再生核且 \texorpdfstring{$\det G$}{det G} 为评价函子体积]\label{prop:xi-hellinger-kernel-rkhs-volume}
沿用命题 \ref{prop:xi-hellinger-kernel-fourier-sech2} 的记号，并记谱密度
\[
w(\xi):=\widehat A(\xi)=\frac{\pi^2}{\cosh^2(\pi\xi)}.
\]
定义 Hilbert 空间
\[
\mathcal H_{A}
:=\Bigl\{f\in\mathcal S'(\RR):\ \|f\|_{\mathcal H_A}^{2}
:=\frac{1}{2\pi}\int_{\RR}\frac{|\widehat f(\xi)|^{2}}{w(\xi)}\,\dd\xi
=\frac{1}{2\pi}\int_{\RR}|\widehat f(\xi)|^{2}\frac{\cosh^{2}(\pi\xi)}{\pi^{2}}\,\dd\xi<\infty\Bigr\}.
\]
则对任意 \(x\in\RR\)，评价泛函 \(f\mapsto f(x)\) 在 \(\mathcal H_A\) 上连续。
令 \(k_x(\cdot):=A(\cdot-x)\)，则 \(k_x\in\mathcal H_A\) 且对所有 \(f\in\mathcal H_A\) 有再生性质
\[
\boxed{\;f(x)=\langle f,k_x\rangle_{\mathcal H_A}\;},
\qquad
\langle k_x,k_y\rangle_{\mathcal H_A}=A(x-y).
\]
给定深度对数谱 \(\boldsymbol{\sigma}=(\sigma_1,\dots,\sigma_\kappa)\)，定义 Gram 矩阵
\[
G(\boldsymbol{\sigma})_{jk}:=A(\sigma_j-\sigma_k)\qquad(1\le j,k\le\kappa).
\]
则
\[
\boxed{\;
\det G(\boldsymbol{\sigma})
=\mathrm{Vol}^{2}\bigl(\mathrm{span}\{k_{\sigma_1},\dots,k_{\sigma_\kappa}\}\bigr),\qquad
S_{1/2}(\boldsymbol{\sigma}):=-\log\det G(\boldsymbol{\sigma}).\;}
\]
\end{proposition}
\begin{proof}
在 \(\mathcal H_A\) 上定义内积
\(
\langle f,g\rangle_{\mathcal H_A}:=\frac{1}{2\pi}\int \widehat f(\xi)\overline{\widehat g(\xi)}\,w(\xi)^{-1}\dd\xi
\)。
对 \(k_x(\cdot)=A(\cdot-x)\)，由 Fourier 平移得 \(\widehat{k_x}(\xi)=w(\xi)e^{-\mathrm{i}\xi x}\)，从而 \(k_x\in\mathcal H_A\) 且
\[
\langle f,k_x\rangle_{\mathcal H_A}
=\frac{1}{2\pi}\int_{\RR}\widehat f(\xi)\,e^{\mathrm{i}\xi x}\,\dd\xi
=f(x),
\]
其中最后一步为 Schwartz 分布的 Fourier 反演公式。
取 \(f=k_x\) 得 \(\langle k_x,k_y\rangle_{\mathcal H_A}=A(x-y)\)。
最后一式为 Hilbert 空间中 Gram 行列式的几何解释：\(\det(\langle v_j,v_k\rangle)=\mathrm{Vol}^{2}(\mathrm{span}\{v_j\})\)。证毕。
\end{proof}

% (User input window)
%
% 本文件把 Gram 行列式体积解释推进到“观测收缩/相干能量/最近邻封口”三条可审计律。

\begin{proposition}[线性观测的体积单调律：Gram--det 熵在收缩算子下严格增大]\label{prop:xi-gram-det-entropy-monotone-under-contraction}
设 \(H\) 为复 Hilbert 空间，\(\psi_1,\dots,\psi_\kappa\in H\) 线性无关，记 Gram 矩阵
\[
G=(\langle\psi_i,\psi_j\rangle)_{1\le i,j\le\kappa}\succ0,
\qquad
S:=-\log\det G.
\]
对任意有界线性算子 \(T:H\to H\) 满足 \(T^\ast T\preceq I\)，令
\[
G_T:=(\langle T\psi_i,T\psi_j\rangle)_{1\le i,j\le\kappa}.
\]
则有
\[
\boxed{\;
0\prec G_T\preceq G,\qquad
\det G_T\le \det G,\qquad
S_T:=-\log\det G_T\ge S.\;}
\]
并且等号成立当且仅当 \(T\) 在子空间 \(\mathrm{span}\{\psi_1,\dots,\psi_\kappa\}\) 上为等距同构（等价于 \(T^\ast T=I\) 在该子空间上的限制）。
因此“带宽截断/噪声平滑/子空间投影”等任何线性收缩观测都统一受制于一条严格的体积不可逆律。
\end{proposition}
\begin{proof}
对任意 \(c=(c_1,\dots,c_\kappa)\in\CC^\kappa\)，置 \(v=\sum_{j=1}^\kappa c_j\psi_j\)。
则
\[
c^\ast G_T c=\|Tv\|^2\le \|v\|^2=c^\ast G c,
\]
故 \(G_T\preceq G\)，且因 \(\{\psi_j\}\) 线性无关得 \(G_T\succ0\)。
由 Loewner 序下特征值单调性（从而行列式单调性）得到 \(\det G_T\le \det G\)，于是 \(S_T\ge S\)。
若 \(\det G_T=\det G\)，则 \(G_T\preceq G\) 且行列式相等迫使 \(G_T=G\)，从而 \(\|Tv\|=\|v\|\) 对所有 \(v\in\mathrm{span}\{\psi_j\}\) 成立，即 \(T\) 在该子空间上为等距同构。证毕。
\end{proof}
\leanverified{paper\_xi\_gram\_det\_entropy\_monotone\_under\_contraction}

\begin{proposition}[相干能量下界：体积熵由平方相干强制]\label{prop:xi-gram-entropy-lower-bound-by-coherence}
沿用命题 \ref{prop:xi-gram-det-entropy-monotone-under-contraction} 的记号，并进一步假设 \(\|\psi_j\|=1\)（\(1\le j\le\kappa\)）。
于是 \(G=I+E\) 且 \(E\) Hermitian、\(\mathrm{diag}(E)=0\)。
则有严格下界
\[
\boxed{\;
S=-\log\det(I+E)\ \ge\ \frac{\|E\|_{\mathrm F}^{2}}{2\,(1+\|E\|_{2})}\;}
\]
其中 \(\|E\|_{\mathrm F}^{2}=\sum_{i\neq j}|\langle\psi_i,\psi_j\rangle|^2\)、\(\|E\|_2\) 为谱范数。
特别地，当 \(\|E\|_2=o(1)\) 时，
\[
S=\frac12\sum_{i\neq j}|\langle\psi_i,\psi_j\rangle|^2+o\!\Bigl(\sum_{i\neq j}|\langle\psi_i,\psi_j\rangle|^2\Bigr),
\]
从而体积熵在“弱相干相”中由二体平方相干完全主导。
\end{proposition}
\begin{proof}[证明要点]
令 \(\{\lambda_r\}_{r=1}^\kappa\) 为 \(E\) 的实特征值，则 \(\lambda_r>-1\)（因 \(I+E\succ0\)），且 \(\sum_r\lambda_r=\mathrm{tr}(E)=0\)。
因此
\[
S=-\sum_{r=1}^\kappa\log(1+\lambda_r)=\sum_{r=1}^\kappa\bigl(\lambda_r-\log(1+\lambda_r)\bigr).
\]
对任意 \(x>-1\) 有标量不等式
\[
x-\log(1+x)\ge \frac{x^2}{2(1+|x|)}\ge \frac{x^2}{2(1+\|E\|_2)}.
\]
求和并用 \(\sum_r\lambda_r^2=\|E\|_{\mathrm F}^2\) 得所示下界。
弱相干展开由矩阵对数级数
\(
\log\det(I+E)=\mathrm{tr}(E)-\frac12\mathrm{tr}(E^2)+O(\|E\|_2\|E\|_{\mathrm F}^2)
\)
与 \(\mathrm{tr}(E)=0\) 得到。证毕。
\end{proof}
\leanverified{paper\_xi\_gram\_entropy\_lower\_bound\_by\_coherence}

\begin{proposition}[最近邻乘积上界：体积熵由相邻二体退化严格封口]\label{prop:xi-gram-det-nearest-neighbor-product-upper-bound}
\leanverified{paper\_xi\_gram\_det\_nearest\_neighbor\_product\_upper\_bound}
设 \(\psi_1,\dots,\psi_\kappa\in H\) 满足 \(\|\psi_j\|=1\)，其 Gram 矩阵为 \(G=(\langle\psi_i,\psi_j\rangle)\)。
则有严格不等式
\[
\boxed{\;
\det G\ \le\ \prod_{j=1}^{\kappa-1}\bigl(1-|\langle\psi_{j+1},\psi_j\rangle|^2\bigr)\;}
\qquad
\Longrightarrow\qquad
S=-\log\det G\ \ge\ \sum_{j=1}^{\kappa-1}-\log\bigl(1-|\langle\psi_{j+1},\psi_j\rangle|^2\bigr).
\]
在本文的 Hellinger--RKHS 设定下取 \(\psi_j=k_{s_j}\) 且 \(s_1<\cdots<s_\kappa\)，则 \(|\langle\psi_{j+1},\psi_j\rangle|=A(\Delta_j)\)（\(\Delta_j=s_{j+1}-s_j\)），从而体积熵势在相邻间隔上被一维最近邻账本严格封口；任一间隔塌缩（\(A(\Delta_j)\uparrow 1\)）即强制右端对数发散，给出簇化病态的最小可审计证书。
\end{proposition}
\begin{proof}
将 \(\{\psi_j\}\) 依次作 Gram--Schmidt 正交化：令
\(
v_1:=\psi_1
\)，以及对 \(j\ge 2\),
\(
v_j:=\psi_j-\mathrm{Proj}_{\mathrm{span}\{\psi_1,\dots,\psi_{j-1}\}}\psi_j
\)。
则 Gram 行列式的体积解释给出
\[
\det G=\prod_{j=1}^{\kappa}\|v_j\|^2.
\]
并且对 \(j\ge 2\),
\[
\|v_j\|^2=1-\bigl\|\mathrm{Proj}_{\mathrm{span}\{\psi_1,\dots,\psi_{j-1}\}}\psi_j\bigr\|^2
\le 1-|\langle\psi_j,\psi_{j-1}\rangle|^2,
\]
因为投影到 \(\mathrm{span}\{\psi_1,\dots,\psi_{j-1}\}\) 的范数至少不小于沿单位向量 \(\psi_{j-1}\) 的分量。
相乘即得所示不等式；取负对数得到熵下界。证毕。
\end{proof}


\begin{proposition}[深度体积熵势的精确一阶变分与平衡方程]\label{prop:xi-hellinger-gram-first-variation-equilibrium}
令深度对数谱 \(\mathbf s=(s_1,\dots,s_\kappa)\) 两两不同，并令 Hellinger--Gram 矩阵
\[
G(\mathbf s)_{jk}=A(s_j-s_k)\qquad(1\le j,k\le \kappa),
\]
其中核 \(A\) 见命题 \ref{prop:xi-hellinger-kernel-fourier-sech2}。
定义深度体积熵势
\[
S_{1/2}(\mathbf s):=-\log\det G(\mathbf s).
\]
则对每个 \(1\le j\le \kappa\) 有严格梯度公式
\[
\boxed{\;
\partial_{s_j}S_{1/2}(\mathbf s)
=-2\sum_{k=1}^{\kappa}\bigl(G(\mathbf s)^{-1}\bigr)_{jk}\,A'(s_j-s_k)\;}
\]
其中 \(A'(0):=0\)，且对 \(\Delta\neq 0\),
\[
\boxed{\;
A'(\Delta)
=A(\Delta)\Bigl(\frac{1}{\Delta}-\frac12\coth(\Delta/2)\Bigr)\; }.
\]
由于 \(S_{1/2}\) 对整体平移 \(\mathbf s\mapsto \mathbf s+c\mathbf 1\) 不变，任一在平移自由度被固定后的驻点必须满足平衡方程
\[
\sum_{k=1}^{\kappa}\bigl(G(\mathbf s)^{-1}\bigr)_{jk}\,A'(s_j-s_k)=0,\qquad j=1,\dots,\kappa.
\]
并且力核 \(\frac{1}{\Delta}-\frac12\coth(\Delta/2)\) 在 \(|\Delta|\to\infty\) 与 \(\Delta\to 0\) 的两端满足
\[
\frac{1}{\Delta}-\frac12\coth(\Delta/2)= -\frac12\,\mathrm{sgn}(\Delta)+O(|\Delta|^{-1}),
\qquad
\frac{1}{\Delta}-\frac12\coth(\Delta/2)=-\frac{\Delta}{12}+O(\Delta^3).
\]
\end{proposition}
\begin{proof}
由矩阵微分恒等式 \(\partial\log\det G=\mathrm{tr}(G^{-1}\partial G)\)，得
\[
\partial_{s_j}\log\det G(\mathbf s)=\mathrm{tr}\bigl(G(\mathbf s)^{-1}\,\partial_{s_j}G(\mathbf s)\bigr).
\]
注意 \(\partial_{s_j}G_{ab}=\delta_{aj}A'(s_j-s_b)-\delta_{bj}A'(s_a-s_j)\)，且 \(A'\) 为奇函数。
于是
\[
\partial_{s_j}\log\det G(\mathbf s)
=\sum_{k=1}^{\kappa}\bigl(G^{-1}\bigr)_{jk}A'(s_j-s_k)
+\sum_{k=1}^{\kappa}\bigl(G^{-1}\bigr)_{kj}A'(s_j-s_k)
=2\sum_{k=1}^{\kappa}\bigl(G^{-1}\bigr)_{jk}A'(s_j-s_k),
\]
从而得到所示 \(\partial_{s_j}S_{1/2}=-\partial_{s_j}\log\det G\) 公式。
对 \(\Delta\neq 0\)，直接对 \(A(\Delta)=\Delta/(2\sinh(\Delta/2))\) 求导并整理得
\(
A'(\Delta)=A(\Delta)\bigl(\Delta^{-1}-\tfrac12\coth(\Delta/2)\bigr)
\)；再由奇对称令 \(A'(0)=0\)。
\par
平衡方程来自驻点条件与平移不变性（等价于在约束 \(\sum_j s_j=0\) 或固定任一基点的条件下取驻点）。
两端展开由 \(\coth x=1+O(e^{-2x})\)（\(x\to\infty\)）与
\(\coth x=x^{-1}+x/3+O(x^3)\)（\(x\to 0\)）代入即得。证毕。
\end{proof}

\begin{remark}[两体势的面积律与对数修正：\texorpdfstring{$1/4$}{1/4} 的长程再现]\label{rem:xi-hellinger-two-body-potential-area-log}
令两体势 \(V(\Delta):=-\log A(\Delta)\)（\(\Delta\ge 0\)）。
则对 \(\Delta>0\) 有严格恒等式
\[
\boxed{\;
V(\Delta)=\frac{\Delta}{2}-\log\Delta+\log\bigl(1-e^{-\Delta}\bigr)\;}
\]
并因此
\[
V(\Delta)=\frac{\Delta}{2}-\log\Delta+O(e^{-\Delta})\qquad(\Delta\to\infty).
\]
若改用物理距离变量 \(d:=2\Delta\)，则
\[
V(d/2)=\frac{d}{4}-\log d+\log 2+O(e^{-d/2})\qquad(d\to\infty),
\]
其中线性项系数 \(1/4\) 在该核的长程势中以严格形式出现，而对数项给出不可消去的次级校正。
\end{remark}

\begin{proposition}[深度能量泛函的严格凸性：\texorpdfstring{$-\log A$}{-log A} 的条件负定性]\label{prop:xi-hellinger-energy-functional-strict-convexity}
定义连续能量泛函
\[
\mathcal E(\nu):=\frac12\iint_{\RR\times\RR} -\log A(x-y)\,\dd\nu(x)\dd\nu(y),
\]
其中 \(\nu\) 为具有有限能量的概率测度。
则核 \(\psi(\Delta):=-\log A(\Delta)\) 为（严格）条件负定函数：对任意实系数 \(c_1,\dots,c_m\) 满足 \(\sum_i c_i=0\) 与任意点 \(x_1,\dots,x_m\in\RR\) 有
\[
\sum_{i,j=1}^{m}c_ic_j\,\psi(x_i-x_j)\le 0,
\]
并因此 \(\mathcal E\) 在概率测度的凸集合上为（严格）凸泛函。
在消除整体平移自由度的任一仿射约束下（例如固定均值），\(\mathcal E\) 的极小测度唯一。
\end{proposition}
\begin{proof}[证明要点]
由 Weierstrass 乘积
\(
2\sinh(\Delta/2)=\Delta\prod_{n\ge 1}\bigl(1+\Delta^2/(4\pi^2n^2)\bigr)
\)
得
\[
A(\Delta)=\prod_{n\ge 1}\Bigl(1+\frac{\Delta^2}{4\pi^2n^2}\Bigr)^{-1}.
\]
每个因子 \((1+b^2\Delta^2)^{-1}\) 是中心 Laplace 分布的特征函数，故 \(A\) 为一族独立 Laplace 变量之和的特征函数，从而为无穷可分；由 L\'evy--Khintchine 理论，\(\psi=-\log A\) 为连续负定函数（等价于所述条件负定不等式）。
\par
对 \(\nu_t=(1-t)\nu_0+t\nu_1\) 展开 \(\mathcal E(\nu_t)\) 可得
\(
\mathcal E(\nu_t)=(1-t)\mathcal E(\nu_0)+t\mathcal E(\nu_1)-\frac{t(1-t)}{2}\iint \psi\,\dd(\nu_1-\nu_0)\dd(\nu_1-\nu_0)
\)；
条件负定性使最后一项非负，从而得到凸性。严格性与约束下的唯一性由等号情形强制 \(\nu_0=\nu_1\)（或仅差整体平移）推出。证毕。
\end{proof}

\begin{theorem}[深度 Hellinger--Gram 的严格全正性与谱振荡律]\label{thm:xi-hellinger-gram-strict-total-positivity}
\leanverified{paper\_xi\_hellinger\_gram\_strict\_total\_positivity}
对严格递增深度 \(s_1<\cdots<s_\kappa\)，矩阵 \(G(\mathbf s)=(A(s_j-s_k))_{1\le j,k\le\kappa}\) 为严格全正矩阵。
因此：
\begin{itemize}
\item 全部特征值均为单根；
\item 第 \(r\) 大特征值对应的特征向量在索引序上恰有 \(r-1\) 次严格符号变号；
\item 任意非零向量经 \(G(\mathbf s)\) 作用后的符号变号数严格不增（变差递减）。
\end{itemize}
\end{theorem}
\begin{proof}[证明要点]
由命题 \ref{prop:xi-hellinger-kernel-fourier-sech2}，归一化特征函数满足
\[
\frac{\widehat A(\xi)}{\widehat A(0)}=\mathrm{sech}^{2}(\pi\xi)
=\prod_{n\ge 1}\Bigl(1+\frac{\xi^{2}}{(n-\tfrac12)^{2}}\Bigr)^{-2},
\]
其中使用 \(\cosh(\pi z)=\prod_{n\ge 1}(1+z^{2}/(n-\tfrac12)^{2})\) 的 Weierstrass 乘积。
因此 \(A/\widehat A(0)\) 为 PF\(_\infty\)（P\'olya frequency）密度；由 Schoenberg 的 PF\(_\infty\)\(\Rightarrow\)严格全正平移核定理可知，核 \((x,y)\mapsto A(x-y)\) 为严格全正。
令 \((x_i)=(y_i)=(s_i)\) 即得 \(G(\mathbf s)\) 的严格全正性。
谱单纯性、特征向量振荡律与变差递减性质为全正矩阵的基本定理（Gantmacher--Kre\u{\i}n 变差递减/振荡定理）的直接推论。证毕。
\end{proof}

\begin{corollary}[逆矩阵的棋盘符号律：交替耦合的代数证书]\label{cor:xi-hellinger-gram-inverse-checkerboard}
在定理 \ref{thm:xi-hellinger-gram-strict-total-positivity} 的条件下，\(G(\mathbf s)^{-1}\) 具有严格棋盘符号：
\[
\boxed{\;
(-1)^{i+j}\bigl(G(\mathbf s)^{-1}\bigr)_{ij}>0\qquad(1\le i,j\le \kappa).\;}
\]
\end{corollary}
\begin{proof}
这是全正矩阵逆的符号正规性（sign-regular）结论：非奇异全正矩阵的逆为严格棋盘符号矩阵。证毕。
\end{proof}

\begin{proposition}[Andreief 表示与频域对数气体：\texorpdfstring{$S_{1/2}$}{S_{1/2}} 为外势自由能]\label{prop:xi-hellinger-kernel-andreief-loggas}
在命题 \ref{prop:xi-hellinger-kernel-rkhs-volume} 的设定下，有严格积分表示
\[
\boxed{\;
\det G(\boldsymbol{\sigma})
=\frac{1}{\kappa!(2\pi)^{\kappa}}
\int_{\RR^\kappa}
\Bigl(\prod_{r=1}^{\kappa} w(\xi_r)\Bigr)\,
\bigl|\det(e^{-\mathrm{i}\xi_r\sigma_j})_{1\le r,j\le\kappa}\bigr|^{2}
\ \dd\xi_1\cdots \dd\xi_\kappa.\;}
\]
因此 \(S_{1/2}(\boldsymbol{\sigma})\) 等价于一维 \(\beta=2\) 对数相互作用体系在外势
\(
V(\xi):=-\log w(\xi)=2\log\cosh(\pi\xi)-\log\pi^{2}
\)
下的配分函数自由能；深度谱 \(\boldsymbol{\sigma}\) 以外部相位插值进入 Vandermonde 型体积因子。
\end{proposition}
\begin{proof}
由命题 \ref{prop:xi-hellinger-kernel-fourier-sech2} 的 Fourier 表示，
\[
A(x-y)=\frac{1}{2\pi}\int_{\RR}w(\xi)\,e^{\mathrm{i}\xi(x-y)}\,\dd\xi.
\]
令测度 \(\dd\mu(\xi):=w(\xi)\dd\xi/(2\pi)\)，并置
\(
F_j(\xi):=e^{-\mathrm{i}\xi\sigma_j}
\)。
则 \(G_{jk}=\int F_j(\xi)\overline{F_k(\xi)}\,\dd\mu(\xi)\)。
由 Andreief 恒等式（亦即 Gram 行列式的积分化）得到
\[
\det(G_{jk})
=\frac{1}{\kappa!}\int_{\RR^\kappa}\bigl|\det(F_j(\xi_r))_{1\le r,j\le\kappa}\bigr|^{2}\,\dd\mu(\xi_1)\cdots\dd\mu(\xi_\kappa),
\]
化为所示形式。自由能表述由对数取负得到。证毕。
\end{proof}

\begin{theorem}[碰撞常数的 Hankel 封口：簇化缩放下 \texorpdfstring{$\det G$}{det G} 的 confluent Vandermonde 因子化]\label{thm:xi-hellinger-gram-collision-hankel-constant}
沿用命题 \ref{prop:xi-hellinger-kernel-andreief-loggas} 的记号，并取任意 \(\kappa\ge 1\)。
给定两两不同的 \(t_1,\dots,t_\kappa\in\RR\) 与任意 \(s_0\in\RR\)，考虑簇化缩放
\[
\sigma_j=s_0+\varepsilon t_j\qquad(1\le j\le \kappa),\qquad \varepsilon\downarrow 0,
\]
并令 \(G(\boldsymbol{\sigma})=(A(\sigma_j-\sigma_k))_{1\le j,k\le\kappa}\)。
则存在严格渐近
\[
\boxed{\;
\det G(\boldsymbol{\sigma})
=\varepsilon^{\kappa(\kappa-1)}
\Bigl(\prod_{j=0}^{\kappa-1}\frac{1}{(j!)^{2}}\Bigr)\,
\mathfrak H_\kappa\,
\prod_{1\le j<k\le\kappa}(t_k-t_j)^{2}\,\bigl(1+o(1)\bigr)\;}
\]
其中普适常数 \(\mathfrak H_\kappa\) 为权函数 \(w(\xi)=\pi^{2}\mathrm{sech}^{2}(\pi\xi)\) 的正规化 Hankel 行列式：
\[
\mathfrak H_\kappa:=\det(\mu_{p+q})_{p,q=0}^{\kappa-1},\qquad
\mu_n:=\frac{1}{2\pi}\int_{\RR}\xi^{n}w(\xi)\,\dd\xi.
\]
因此碰撞发散的 Vandermonde 幂律之外，仍存在仅依赖于 \(\kappa\) 的不可约体积常数 \(\mathfrak H_\kappa\)。
\end{theorem}
\begin{proof}
由命题 \ref{prop:xi-hellinger-kernel-andreief-loggas} 的 Andreief 表示与平移不变性，\(\det G(\boldsymbol{\sigma})\) 仅依赖差分 \(\sigma_j-\sigma_k=\varepsilon(t_j-t_k)\)，且
\[
\det G(\boldsymbol{\sigma})
=\frac{1}{\kappa!(2\pi)^{\kappa}}
\int_{\RR^\kappa}\Bigl(\prod_{r=1}^{\kappa}w(\xi_r)\Bigr)\,
\bigl|\det(e^{-\mathrm{i}\varepsilon\xi_r t_j})_{1\le r,j\le\kappa}\bigr|^{2}\,\dd\xi_1\cdots\dd\xi_\kappa.
\]
当 \(\varepsilon\downarrow 0\) 时，经典的 confluent Vandermonde 渐近给出
\[
\det(e^{-\mathrm{i}\varepsilon\xi_r t_j})_{r,j}
=\varepsilon^{\kappa(\kappa-1)/2}\Bigl(\prod_{j=0}^{\kappa-1}\frac{1}{j!}\Bigr)
\Delta(\boldsymbol{\xi})\,\Delta(\mathbf t)\,\bigl(1+o(1)\bigr),
\]
其中 \(\Delta(\boldsymbol{\xi})=\prod_{1\le r<\ell\le\kappa}(\xi_\ell-\xi_r)\)、\(\Delta(\mathbf t)=\prod_{1\le j<k\le\kappa}(t_k-t_j)\)。
取模平方并代回，得到
\[
\det G(\boldsymbol{\sigma})
=\varepsilon^{\kappa(\kappa-1)}\Bigl(\prod_{j=0}^{\kappa-1}\frac{1}{(j!)^{2}}\Bigr)\,
\Delta(\mathbf t)^{2}\cdot
\frac{1}{\kappa!}\int_{\RR^\kappa}\Delta(\boldsymbol{\xi})^{2}\,\dd\mu(\xi_1)\cdots\dd\mu(\xi_\kappa)\,\bigl(1+o(1)\bigr),
\]
其中 \(\dd\mu(\xi)=w(\xi)\dd\xi/(2\pi)\)。
最后，由 Andreief 恒等式（Vandermonde 平方的矩阵化积分）有
\[
\frac{1}{\kappa!}\int_{\RR^\kappa}\Delta(\boldsymbol{\xi})^{2}\,\dd\mu(\xi_1)\cdots\dd\mu(\xi_\kappa)
=\det(\mu_{p+q})_{p,q=0}^{\kappa-1}=\mathfrak H_\kappa,
\]
从而得到所示渐近。证毕。
\end{proof}

\begin{corollary}[簇化相的最小奇异指数：条件数的制度性幂次爆炸]\label{cor:xi-hellinger-gram-cluster-conditioning-exponent}
在定理 \ref{thm:xi-hellinger-gram-collision-hankel-constant} 的簇化缩放下，
存在常数 \(c(\mathbf t)>0\) 使
\[
\det G(\boldsymbol{\sigma})=c(\mathbf t)\,\varepsilon^{\kappa(\kappa-1)}\bigl(1+o(1)\bigr).
\]
并且由于 \(\mathrm{tr}(G)=\kappa\)，最大特征值满足 \(\lambda_{\max}(G)\le \kappa\)。
因此最小特征值具有硬上界
\[
\lambda_{\min}(G)\le \frac{\det G}{\lambda_{\max}(G)^{\kappa-1}}
\le \kappa^{1-\kappa}\,c(\mathbf t)\,\varepsilon^{\kappa(\kappa-1)}\bigl(1+o(1)\bigr),
\]
从而谱条件数满足
\[
\boxed{\;
\mathrm{cond}_2(G)=\frac{\lambda_{\max}(G)}{\lambda_{\min}(G)}
\ \ge\ C(\mathbf t)\,\varepsilon^{-\kappa(\kappa-1)}\bigl(1+o(1)\bigr)\;}
\qquad(\varepsilon\downarrow 0),
\]
其中 \(C(\mathbf t)>0\) 为常数。
因此簇化相的数值不适定性由指数 \(\kappa(\kappa-1)\) 严格锁定。
\end{corollary}
\begin{proof}
行列式与谱乘积给出
\(
\det G=\lambda_{\min}\prod_{r\neq\min}\lambda_r\ge \lambda_{\min}\lambda_{\max}^{\kappa-1}
\)，
故 \(\lambda_{\min}\le \det G/\lambda_{\max}^{\kappa-1}\)。
又 \(\lambda_{\max}\le \mathrm{tr}(G)=\kappa\)，并代入定理 \ref{thm:xi-hellinger-gram-collision-hankel-constant} 的渐近，即得所示条件数下界。证毕。
\end{proof}

\begin{proposition}[普适 Hankel 常数的算术性：\texorpdfstring{$\mathfrak H_\kappa\in\QQ$}{H_k in Q} 且由偶 \texorpdfstring{$\zeta$}{zeta} 值显式封口]\label{prop:xi-hellinger-hankel-constant-rationality}
\leanverified{paper\_xi\_hellinger\_hankel\_constant\_rationality}
沿用定理 \ref{thm:xi-hellinger-gram-collision-hankel-constant} 的记号。
则矩 \(\mu_n\) 满足 \(\mu_{2n+1}=0\)，且对任意 \(n\ge 0\) 有严格闭式
\[
\boxed{\;
\mu_{2n}
=2^{1-2n}\pi^{-2n}(2n)!\,\bigl(1-2^{1-2n}\bigr)\zeta(2n)\in\QQ\; }.
\]
从而对任意 \(\kappa\ge 1\)，\(\mathfrak H_\kappa=\det(\mu_{p+q})_{p,q=0}^{\kappa-1}\in\QQ\)。
例如
\[
\boxed{\;
\mathfrak H_1=1,\quad
\mathfrak H_2=\frac{1}{12},\quad
\mathfrak H_3=\frac{1}{540},\quad
\mathfrak H_4=\frac{1}{42000},\quad
\mathfrak H_5=\frac{1}{3215625}.\;}
\]
\end{proposition}
\begin{proof}[证明要点]
奇次矩为零来自 \(w(\xi)\) 的偶对称。
对偶次矩，可由 \(\mathrm{sech}^2\) 的部分分式/留数展开（等价地，由 \(\tanh\) 的晶格极点结构）将
\(\int_{\RR}\xi^{2n}\mathrm{sech}^{2}(\pi\xi)\dd\xi\)
化为 Hurwitz \(\zeta\) 值之差，从而得到所示 \(\zeta(2n)\) 闭式；再由
\(\zeta(2n)\in \pi^{2n}\QQ\)
推出 \(\mu_{2n}\in\QQ\)。
由 Hankel 行列式对矩的多项式性即得 \(\mathfrak H_\kappa\in\QQ\)。
所列数值由该闭式直接计算。证毕。
\end{proof}

\begin{remark}[算术带宽定律：\texorpdfstring{$\mathfrak H_\kappa$}{H_k} 的 \texorpdfstring{$p$}{p}-进刚性仅支持于 \texorpdfstring{$p\le 2\kappa-1$}{p<=2k-1}]\label{rem:xi-hellinger-hankel-padic-bandwidth}
令 \(\nu_p(\cdot)\) 为 \(p\)-进赋值。
由命题 \ref{prop:xi-hellinger-hankel-constant-rationality} 中 \(\mu_{2n}\) 的 Bernoulli--von Staudt--Clausen 分母结构可知：
对任意素数 \(p>2\kappa-1\)，有 \(\nu_p(\mu_0),\dots,\nu_p(\mu_{2\kappa-2})\ge 0\)。
因此 \(\nu_p(\mathfrak H_\kappa)=0\)（\(p>2\kappa-1\)）。
换言之，\(\mathfrak H_\kappa\) 的全部 \(p\)-进“病态/刚性”只能由有限素数集合 \(\{p:\ p\le 2\kappa-1\}\) 贡献。
\end{remark}

\begin{proposition}[等距深度的 Toeplitz 化：\texorpdfstring{$\cosh^{-2}$}{cosh^{-2}} 的 Poisson 周期和符号]\label{prop:xi-hellinger-toeplitz-symbol-poisson}
若 \(\sigma_j=\sigma_0+(j-1)\Delta\)（\(\Delta>0\)），则 \(G(\boldsymbol{\sigma})\) 为 Toeplitz 矩阵
\(
G_{jk}=A((j-k)\Delta)
\)。
其 Toeplitz 符号 \(\sigma_\Delta:[-\pi,\pi]\to(0,\infty)\) 可写为
\[
\boxed{\;
\sigma_\Delta(\theta)
:=\sum_{n\in\ZZ}A(n\Delta)e^{\mathrm{i}n\theta}
=\frac{1}{\Delta}\sum_{k\in\ZZ}\frac{\pi^{2}}{\cosh^{2}\!\Bigl(\pi\frac{\theta+2\pi k}{\Delta}\Bigr)}\;}
\qquad(\theta\in[-\pi,\pi]).
\]
等价地，令 \(q:=e^{-\Delta/2}\) 并取 Jacobi \(\vartheta_4\) 函数（nome \(q\)）
\[
\vartheta_4(z,q):=\sum_{m\in\ZZ}(-1)^m q^{m^2}e^{2\mathrm{i}mz},
\]
则对所有 \(\theta\in[-\pi,\pi]\) 有严格恒等式
\[
\boxed{\;
\sigma_\Delta(\theta)
=1+\Delta\,\partial_\theta^{2}\log\vartheta_4\!\Bigl(\frac{\theta}{2},q\Bigr)
=1+2\Delta\sum_{n\ge 1}\frac{nq^{n}}{1-q^{2n}}\cos(n\theta).\;}
\]
\end{proposition}
\begin{proof}
Toeplitz 性由 \(G_{jk}=A(\sigma_j-\sigma_k)=A((j-k)\Delta)\) 直接得到。
Poisson 求和公式在本 Fourier 约定下给出
\(
\sum_{n\in\ZZ}A(n\Delta)e^{\mathrm{i}n\theta}
=\Delta^{-1}\sum_{k\in\ZZ}\widehat A((\theta+2\pi k)/\Delta)
\)。
代入命题 \ref{prop:xi-hellinger-kernel-fourier-sech2} 的 \(\widehat A=w\) 即得 Poisson 周期和表示。
\par
另一方面，由
\(
A(n\Delta)=\frac{n\Delta}{2\sinh(n\Delta/2)}
=n\Delta\,\frac{q^n}{1-q^{2n}}
\)
（\(n\ge 1\)、\(q=e^{-\Delta/2}\)）得 Fourier 级数表示
\[
\sigma_\Delta(\theta)=1+2\sum_{n\ge 1}A(n\Delta)\cos(n\theta)
=1+2\Delta\sum_{n\ge 1}\frac{nq^{n}}{1-q^{2n}}\cos(n\theta).
\]
再由 \(\vartheta_4\) 的 Jacobi 乘积公式推出标准恒等式
\[
\partial_z^{2}\log\vartheta_4(z,q)
=8\sum_{n\ge 1}\frac{nq^{n}}{1-q^{2n}}\cos(2nz),
\]
取 \(z=\theta/2\) 并用 \(\partial_\theta^{2}=(1/4)\partial_z^{2}\) 即得导数闭式。证毕。
\end{proof}

\begin{remark}[椭圆延拓与半周期二阶极点晶格]\label{rem:xi-hellinger-toeplitz-elliptic-extension-halfperiod-poles}
把 \(\sigma_\Delta(\theta)\) 视为复变量 \(\theta\) 的函数。
由 \(\vartheta_4\) 的准周期性可知 \(\partial_\theta^2\log\vartheta_4(\theta/2,q)\) 消去仿射漂移项，从而 \(\sigma_\Delta\) 具有基本周期 \(2\pi\) 与 \(\mathrm{i}\Delta\)。
记矩形晶格
\[
\Lambda_\Delta:=2\pi\ZZ+\mathrm{i}\Delta\ZZ,
\qquad
E_\Delta:=\CC/\Lambda_\Delta.
\]
则 \(\sigma_\Delta\) 在椭圆曲线 \(E_\Delta\) 上诱导一个亚纯函数；
并且由命题 \ref{prop:xi-hellinger-toeplitz-symbol-poisson} 的 Poisson 周期和表示可知其极点均为二阶极点且恰位于半周期平移晶格
\[
\theta\in \mathrm{i}\Delta\Bigl(\ZZ+\tfrac12\Bigr)+2\pi\ZZ.
\]
\end{remark}

\begin{remark}[模对偶尺度 \texorpdfstring{$\Delta\leftrightarrow 2\pi^{2}/\Delta$}{Delta <-> 2pi^2/Delta}：Poisson 尖峰与 Fourier 衰减的强制等价]\label{rem:xi-hellinger-toeplitz-modular-duality-scale}
在命题 \ref{prop:xi-hellinger-toeplitz-symbol-poisson} 中令 \(\tau:=\mathrm{i}\Delta/(2\pi)\)，则 nome \(q=e^{\mathrm{i}\pi\tau}=e^{-\Delta/2}\)。
Jacobi 变换 \(\tau\mapsto -1/\tau\) 导出对偶 nome
\[
q_\star:=e^{\mathrm{i}\pi(-1/\tau)}=e^{-2\pi^{2}/\Delta}.
\]
在该变换下，\(\sigma_\Delta\) 的 Fourier 级数表示与 Poisson 周期和表示互为同一 θ-结构的两种等价表象；
其中尺度常数 \(2\pi^{2}/\Delta\) 并非近似技巧，而由 θ-函数的模对偶性制度性地强制出现。
\end{remark}

\begin{proposition}[等距 Toeplitz 相的逆矩阵局域化：反演算子的指数衰减由带宽封口]\label{prop:xi-hellinger-toeplitz-inverse-localization}
在命题 \ref{prop:xi-hellinger-toeplitz-symbol-poisson} 的等距设定下，令
\(
G_\kappa(\Delta)=(A((j-k)\Delta))_{1\le j,k\le \kappa}
\)
并记其符号为 \(\sigma_\Delta\)。
则 \(\sigma_\Delta\) 在带状域 \(|\Im\theta|<\Delta/2\) 内解析。
并且对任意 \(0<\varepsilon<\Delta/2\)，存在常数 \(C_{\Delta,\varepsilon}>0\) 使对一切 \(\kappa\) 与 \(1\le j,k\le\kappa\) 有
\[
\boxed{\;
\bigl|\bigl(G_\kappa(\Delta)^{-1}\bigr)_{jk}\bigr|
\le C_{\Delta,\varepsilon}\,e^{-(\Delta/2-\varepsilon)|j-k|}\;}
\]
从而等距深度编码具有严格的“局域可逆性”。
指数门槛 \(\Delta/2\) 由 \(\sigma_\Delta\) 的最近复奇点距离封口（见命题 \ref{prop:xi-hellinger-toeplitz-symbol-poisson} 的 Fourier 系数衰减率）。
\end{proposition}
\begin{proof}[证明要点]
由 \(A(n\Delta)=O(n\Delta\,e^{-n\Delta/2})\) 得 \(\sigma_\Delta\) 的 Fourier 系数指数衰减，从而 \(\sigma_\Delta\) 在 \(|\Im\theta|<\Delta/2\) 内解析。
并且由于 \(\sigma_\Delta(\theta)>0\)（\(\theta\in[-\pi,\pi]\)）且解析函数连续，任取 \(0<\varepsilon<\Delta/2\)，在闭带 \(|\Im\theta|\le \Delta/2-\varepsilon\) 上有一致下界 \(\inf|\sigma_\Delta|>0\)。
因此 \(1/\sigma_\Delta\) 在该闭带内解析且有界，其 Fourier 系数以速率 \(e^{-(\Delta/2-\varepsilon)|n|}\) 衰减。
对无限 Toeplitz 算子，逆核的 Toeplitz 系数即为 \(1/\sigma_\Delta\) 的 Fourier 系数；有限截断 \(G_\kappa(\Delta)^{-1}\) 的矩阵元衰减由同一 Wiener 代数闭包与有限截断稳定性推出。证毕。
\end{proof}

\begin{proposition}[解析带宽 \texorpdfstring{$\Delta/2$}{Delta/2} 的尖锐性与指数门槛不可超越]\label{prop:xi-hellinger-toeplitz-bandwidth-sharpness}
在命题 \ref{prop:xi-hellinger-toeplitz-symbol-poisson} 的等距设定下，定义解析带宽
\[
a_{\max}:=\sup\Bigl\{a>0:\ \sigma_\Delta\ \text{在}\ |\Im\theta|<a\ \text{内解析}\Bigr\}.
\]
则 \(a_{\max}=\Delta/2\)。
并且不存在 \(\delta>0\) 使得对所有 \(\kappa\) 与所有 \(1\le j,k\le \kappa\) 有一致界
\[
\bigl|\bigl(G_\kappa(\Delta)^{-1}\bigr)_{jk}\bigr|\ \le\ C\,e^{-(\Delta/2+\delta)|j-k|}.
\]
\end{proposition}
\begin{proof}
由 Poisson 周期和表示
\(
\sigma_\Delta(\theta)=\Delta^{-1}\sum_{k\in\ZZ}\pi^2\cosh^{-2}(\pi(\theta+2\pi k)/\Delta)
\)
与 \(\cosh(u)=0\iff u=\mathrm{i}\pi(\tfrac12+n)\)（\(n\in\ZZ\)）可知，\(\sigma_\Delta\) 在
\(
\theta=-2\pi k+\mathrm{i}\Delta(\tfrac12+n)
\)
处具有二阶极点。
因此在任意 \(a>\Delta/2\) 的条带 \(|\Im\theta|<a\) 内必包含极点，\(\sigma_\Delta\) 不可能解析；结合命题 \ref{prop:xi-hellinger-toeplitz-inverse-localization} 中 \(|\Im\theta|<\Delta/2\) 的解析性，得到 \(a_{\max}=\Delta/2\)。
\par
若存在 \(\delta>0\) 与常数 \(C\) 使所示一致界成立，则在 \(\kappa\to\infty\) 的 Toeplitz 有限截断稳定口径下，双边 Toeplitz 卷积算子 \(T_\Delta\) 的逆核系数将以速率 \(e^{-(\Delta/2+\delta)|n|}\) 绝对可和，从而 \(1/\sigma_\Delta\) 属于带权 Wiener 代数
\(
\sum_{n\in\ZZ}\abs{\widehat{(1/\sigma_\Delta)}(n)}\,e^{(\Delta/2+\delta)|n|}<\infty
\)。
由 Wiener 引理，对应的逆元 \(\sigma_\Delta\) 亦必须属于同一带权代数。
但命题 \ref{prop:xi-hellinger-toeplitz-symbol-poisson} 给出 \(\sigma_\Delta(\theta)=\sum_{n\in\ZZ}A(n\Delta)e^{\mathrm{i}n\theta}\)，而 \(A(n\Delta)\sim n\Delta\,e^{-n\Delta/2}\)（\(n\to\infty\)）使得
\(
\sum_{n\ge 1}A(n\Delta)\,e^{(\Delta/2+\delta)n}=+\infty
\)，矛盾。
证毕。
\end{proof}

\begin{theorem}[Toeplitz 热力学极限：等距体积熵密度由单一符号积分封口]\label{thm:xi-hellinger-toeplitz-szego-thermo-limit}
在命题 \ref{prop:xi-hellinger-toeplitz-symbol-poisson} 的设定下，令
\(
D_\kappa(\Delta):=\det(G_{jk})_{1\le j,k\le\kappa}
\)。
则存在极限
\[
\boxed{\;
\lim_{\kappa\to\infty}\frac{1}{\kappa}\log D_\kappa(\Delta)
=\frac{1}{2\pi}\int_{-\pi}^{\pi}\log\sigma_\Delta(\theta)\,\dd\theta\;}
\]
并因此
\[
\boxed{\;
\lim_{\kappa\to\infty}\frac{1}{\kappa}S_{1/2}(\boldsymbol{\sigma})
=-\frac{1}{2\pi}\int_{-\pi}^{\pi}\log\sigma_\Delta(\theta)\,\dd\theta\;}
\qquad(\sigma_j=\sigma_0+(j-1)\Delta).
\]
\end{theorem}
\begin{proof}
由命题 \ref{prop:xi-hellinger-toeplitz-symbol-poisson}，\(\sigma_\Delta\) 在 \([-\pi,\pi]\) 上严格正且解析。
因此 Szeg\H{o} 定理适用，给出 Toeplitz 行列式的线性极限式；第二式由 \(S_{1/2}=-\log D_\kappa\) 直接推出。证毕。
\end{proof}

\begin{proposition}[Heine--Szeg\H{o} 幺正矩阵模型：等距深度体积即 CUE 外势配分函数]\label{prop:xi-hellinger-toeplitz-heine-szego-cue}
在命题 \ref{prop:xi-hellinger-toeplitz-symbol-poisson} 的设定下，有幺正积分表示
\[
\boxed{\;
D_\kappa(\Delta)
=\frac{1}{(2\pi)^\kappa\kappa!}
\int_{[-\pi,\pi]^\kappa}
\Bigl(\prod_{r=1}^{\kappa}\sigma_\Delta(\theta_r)\Bigr)
\prod_{1\le r<\ell\le\kappa}\bigl|e^{\mathrm{i}\theta_r}-e^{\mathrm{i}\theta_\ell}\bigr|^2\,
\dd\theta_1\cdots\dd\theta_\kappa.\;}
\]
从而等距构型下的 \(S_{1/2}\) 等价于 CUE（\(\beta=2\)）模型在外势 \(-\log\sigma_\Delta\) 下的自由能增量。
\end{proposition}
\begin{proof}
这是 Toeplitz 行列式与幺正本征角积分之间的 Heine--Szeg\H{o} 恒等式。证毕。
\end{proof}

% (User input window)
%
% 本文件承载等距 Toeplitz 体积的尾部渐近：大分离簇展开、准正交相阈值与周期化常数。

\begin{proposition}[大分离簇展开：\texorpdfstring{$S_{1/2}$}{S_{1/2}} 的主项为二体指数能量]\label{prop:xi-hellinger-cluster-expansion-large-separation}
\leanverified{paper\_xi\_hellinger\_cluster\_expansion\_large\_separation}
对一般深度谱 \(\boldsymbol{\sigma}=(\sigma_1,\dots,\sigma_\kappa)\)，记最小分离
\(
\Delta_{\min}:=\min_{j\neq k}|\sigma_j-\sigma_k|
\)。
当 \(\Delta_{\min}\to\infty\) 时，存在一致展开
\[
\boxed{\;
S_{1/2}(\boldsymbol{\sigma})
=\frac12\sum_{j\neq k}A(\sigma_j-\sigma_k)^{2}
+O\!\Bigl(\kappa^{3}(1+\Delta_{\min})^{3}e^{-3\Delta_{\min}/2}\Bigr).\;}
\]
并且
\(
A(\Delta)\sim |\Delta|e^{-|\Delta|/2}
\)，从而主项满足
\[
\frac12\sum_{j\neq k}A(\sigma_j-\sigma_k)^2
\sim \frac12\sum_{j\neq k}|\sigma_j-\sigma_k|^2e^{-|\sigma_j-\sigma_k|}
\qquad(\Delta_{\min}\to\infty).
\]
\end{proposition}
\begin{proof}
写 \(G(\boldsymbol{\sigma})=I+E\)，其中 \(E_{jj}=0\)、\(E_{jk}=A(\sigma_j-\sigma_k)\)（\(j\neq k\)）。
由 \(A(\Delta)=|\Delta|/(2\sinh(|\Delta|/2))\) 得上界
\(
0<A(\Delta)\le C(1+|\Delta|)e^{-|\Delta|/2}
\)
（\(|\Delta|\to\infty\) 的双曲函数渐近给出）。
因此 \(\|E\|_{\max}\le C(1+\Delta_{\min})e^{-\Delta_{\min}/2}\)，并在 \(\Delta_{\min}\) 足够大时 \(\|E\|\) 可取小。
由对数行列式展开
\(
-\log\det(I+E)=\sum_{m\ge 1}\frac{(-1)^m}{m}\tr(E^m)
\)
与 \(\tr(E)=0\) 得首项为 \(\tr(E^2)/2=\frac12\sum_{j\neq k}E_{jk}^2\)。
余项由 \(|\tr(E^m)|\le \kappa^m\|E\|_{\max}^m\) 估计，得到所示误差界。证毕。
\end{proof}

\begin{corollary}[正交化阈值与体积饱和相：\texorpdfstring{$\Delta_{\min}\gtrsim 2\log\kappa$}{Delta_min ~ 2 log kappa} 强制 \texorpdfstring{$\det G\to 1$}{det G -> 1}]\label{cor:xi-hellinger-quasi-orthogonal-phase}
沿用命题 \ref{prop:xi-hellinger-cluster-expansion-large-separation} 的记号。
则存在常数 \(C>0\) 使
\[
\max_{1\le j\le \kappa}\sum_{k\neq j}\bigl|G(\boldsymbol{\sigma})_{jk}\bigr|
\le C\,\kappa(1+\Delta_{\min})e^{-\Delta_{\min}/2}.
\]
因此若构型随 \(\kappa\) 变化并满足 \(\kappa(1+\Delta_{\min})e^{-\Delta_{\min}/2}\to 0\)，则
\[
\lambda_{\min}(G(\boldsymbol{\sigma}))\to 1,\qquad
\lambda_{\max}(G(\boldsymbol{\sigma}))\to 1,\qquad
\det G(\boldsymbol{\sigma})\to 1,\qquad
S_{1/2}(\boldsymbol{\sigma})\to 0.
\]
特别地，充分条件
\[
\Delta_{\min}\ge 2\log\kappa+\omega(1)
\]
强制进入“准正交相”（体积饱和相）：原子数的增长一旦被指数尾 \(e^{-\Delta/2}\) 压制，体积熵势即塌缩为零熵极限。
\end{corollary}
\begin{proof}
由 \(G=I+E\) 且 \(E_{jk}=A(\sigma_j-\sigma_k)\)（\(j\neq k\)），以及命题 \ref{prop:xi-hellinger-cluster-expansion-large-separation} 证明中的上界
\(
A(\Delta)\le C(1+|\Delta|)e^{-|\Delta|/2}
\)，
对每个 \(j\) 有
\[
\sum_{k\neq j}|G_{jk}|=\sum_{k\neq j}A(\sigma_j-\sigma_k)\le (\kappa-1)\,C(1+\Delta_{\min})e^{-\Delta_{\min}/2},
\]
得到第一式。
当右端趋于 \(0\) 时，矩阵范数满足 \(\|E\|_2\le \|E\|_\infty\to 0\)，故谱 \(\mathrm{Spec}(G)\subset [1-\|E\|_2,1+\|E\|_2]\)。
由此 \(\lambda_{\min},\lambda_{\max}\to 1\)，并由 \(\det G=\prod_r\lambda_r\) 推出 \(\det G\to 1\)，从而 \(S_{1/2}=-\log\det G\to 0\)。
充分条件由 \(e^{-\Delta_{\min}/2}\le \kappa^{-1}e^{-\omega(1)/2}\) 立得。证毕。
\end{proof}

\begin{remark}[等距符号的跨域分辨率常数：周期化尾项以 \texorpdfstring{$\exp(-2\pi^2/\Delta)$}{exp(-2pi^2/Delta)} 控制]\label{rem:xi-hellinger-toeplitz-periodization-constant-2pi2}
令 \(\sigma_\Delta^{(0)}(\theta):=\frac{\pi^2}{\Delta}\cosh^{-2}(\pi\theta/\Delta)\) 为命题 \ref{prop:xi-hellinger-toeplitz-symbol-poisson} 中 Poisson 周期和的 \(k=0\) 主项。
则对任意 \(\theta\in[-\pi,\pi]\) 有
\[
0\le \sigma_\Delta(\theta)-\sigma_\Delta^{(0)}(\theta)
:=\frac{1}{\Delta}\sum_{k\neq 0}\frac{\pi^{2}}{\cosh^{2}\!\Bigl(\pi\frac{\theta+2\pi k}{\Delta}\Bigr)}
\le C\,\frac{\pi^{2}}{\Delta}\,e^{-2\pi^{2}/\Delta},
\]
其中 \(C>0\) 与 \(\Delta\) 无关。
因此周期化效应的最小跨域常数为 \(2\pi^2\)：其既控制小 \(\Delta\) 时的单胞主导，也给出尾部周期项进入的统一指数门槛。
\par
特别地在 \(\theta=0\) 处，有精确恒等式
\[
\sigma_\Delta(0)=\sum_{n\in\ZZ}A(n\Delta)
=\frac{\pi^{2}}{\Delta}\sum_{k\in\ZZ}\mathrm{sech}^{2}\!\Bigl(\frac{2\pi^{2}k}{\Delta}\Bigr),
\]
从而当 \(\Delta\downarrow 0\) 时
\[
\sigma_\Delta(0)=\frac{\pi^{2}}{\Delta}\Bigl(1+O(e^{-4\pi^{2}/\Delta})\Bigr),
\qquad
\log\sigma_\Delta(0)=\log\frac{\pi^{2}}{\Delta}+O(e^{-4\pi^{2}/\Delta}),
\]
其中跨胞修正由对偶 nome \(q_\star=e^{-2\pi^{2}/\Delta}\) 以 \(O(q_\star^{2})\) 封口。
\end{remark}

\begin{remark}[\texorpdfstring{$\zeta(4)$}{zeta(4)} 的首入侵律：频域外势的四阶系数强制出现]\label{rem:xi-hellinger-logw-zeta4-entry}
令 \(w(\xi)=\pi^2/\cosh^2(\pi\xi)\)，则在 \(\xi=0\) 处有 Taylor 展开
\[
\log w(\xi)=\log\pi^2-\pi^2\xi^2+\frac{\pi^4}{6}\xi^4+O(\xi^6)
=\log\pi^2-\pi^2\xi^2+15\,\zeta(4)\,\xi^4+O(\xi^6).
\]
因此以二次外势（高斯近似）为基底的任何制度，其首个不可消去的系统性修正必出现在四阶，并由 \(\zeta(4)\) 精确封口；
该谱侧四阶首入侵与注 \ref{rem:xi-scale-q4-threshold-principle} 的四阶闭合阈值相容。
\end{remark}

\begin{proposition}[等距深度 Toeplitz 符号的极值与模指数：\texorpdfstring{$\inf\sigma_\Delta$}{inf sigma} 的 \texorpdfstring{$e^{-2\pi^2/\Delta}$}{exp(-2pi^2/Delta)} 刻度]\label{prop:xi-hellinger-toeplitz-symbol-extrema-qstar}
沿用命题 \ref{prop:xi-hellinger-toeplitz-symbol-poisson} 的记号。
则 \(\sigma_\Delta\) 在 \(\theta=0\) 处取全局最大、在 \(\theta=\pi\) 处取全局最小：
\[
\sigma_\Delta(0)=\sup_{\theta\in[-\pi,\pi]}\sigma_\Delta(\theta),
\qquad
\sigma_\Delta(\pi)=\inf_{\theta\in[-\pi,\pi]}\sigma_\Delta(\theta).
\]
并且当 \(\Delta\downarrow 0\) 时（令对偶 nome \(q_\star:=e^{-2\pi^2/\Delta}\)），有一致渐近
\[
\sigma_\Delta(0)=\frac{\pi^2}{\Delta}\Bigl(1+O(q_\star^{2})\Bigr),
\qquad
\sigma_\Delta(\pi)=\frac{8\pi^2}{\Delta}\,q_\star\Bigl(1+O(q_\star^{2})\Bigr).
\]
因此“最小谱高度”由模对偶尺度 \(2\pi^2/\Delta\) 以 \(q_\star\) 的制度指数精确封口。
\end{proposition}
\begin{proof}[证明要点]
\(\theta=0\) 的极大性可由 Fourier 表示
\(
\sigma_\Delta(\theta)=1+2\sum_{n\ge 1}A(n\Delta)\cos(n\theta)
\)
与 \(A(n\Delta)>0\)、\(\cos(n\theta)\le 1\) 直接推出。
\(\theta=\pi\) 的极小性可由 Poisson 表示
\(
\sigma_\Delta(\theta)=\frac{\pi^2}{\Delta}\sum_{k\in\ZZ}\mathrm{sech}^2\!\bigl(\pi(\theta+2\pi k)/\Delta\bigr)
\)
结合 \(\mathrm{sech}^2\) 的严格单峰性与晶格对称性得到：\(\theta=\pi\) 为到 \(2\pi\ZZ\) 的最远点，从而周期化单峰和在该点取最小。
\par
对 \(\theta=\pi\)，由 Poisson 形式得精确恒等式
\[
\sigma_\Delta(\pi)=\frac{\pi^2}{\Delta}\sum_{k\in\ZZ}\mathrm{sech}^2\!\Bigl(\frac{(2k+1)\pi^2}{\Delta}\Bigr).
\]
当 \(\Delta\downarrow 0\) 时，主贡献来自 \(k=0,-1\) 两项；并用
\(
\mathrm{sech}^2(x)=4e^{-2x}(1+O(e^{-2x}))
\)
即可得到
\(
\sigma_\Delta(\pi)=\frac{8\pi^2}{\Delta}e^{-2\pi^2/\Delta}(1+O(e^{-4\pi^2/\Delta}))
\)。
\(\sigma_\Delta(0)\) 的渐近已由注 \ref{rem:xi-hellinger-toeplitz-periodization-constant-2pi2} 给出。证毕。
\end{proof}

\begin{proposition}[等距深度的条件数指数律：可逆性预算被 \texorpdfstring{$2\pi^2/\Delta$}{2pi^2/Delta} 线性锁定]\label{prop:xi-hellinger-toeplitz-conditioning-exponential-law}
令 \(T_\Delta\) 为双边 Toeplitz 卷积算子
\[
(T_\Delta x)_j=\sum_{n\in\ZZ}A((j-n)\Delta)\,x_n,
\]
其符号为 \(\sigma_\Delta\)。
则 \(\mathrm{Spec}(T_\Delta)=\sigma_\Delta([-\pi,\pi])\)，从而条件数为
\[
\mathrm{cond}(T_\Delta)=\frac{\sup_\theta\sigma_\Delta(\theta)}{\inf_\theta\sigma_\Delta(\theta)}
=\frac{\sigma_\Delta(0)}{\sigma_\Delta(\pi)}
\sim \frac18\,e^{2\pi^2/\Delta}\qquad(\Delta\downarrow 0).
\]
对有限维 Toeplitz 主块 \(G_\kappa(\Delta)=(A((j-k)\Delta))_{1\le j,k\le\kappa}\)，其极端特征值在 \(\kappa\to\infty\) 下收敛到 \(\sup\sigma_\Delta,\inf\sigma_\Delta\)，故在大 \(\kappa\) 下 \(\mathrm{cond}(G_\kappa(\Delta))\) 继承同一指数标度。
于是任何要求相对误差 \(\le\varepsilon\) 的稳定反演都必满足精度位数下界
\[
B\ \gtrsim\ \frac{2\pi^2}{\Delta\ln 2}\ +\ O\bigl(1+\log(1/\varepsilon)\bigr),
\]
其中首项斜率 \(2\pi^2/\ln2\) 为不可改写常数。
\end{proposition}
\begin{proof}[证明要点]
谱结论为 Toeplitz 算子的标准事实：\(T_\Delta\) 在 Fourier 侧为乘子 \(\sigma_\Delta\)，故谱为其值域。
条件数渐近由命题 \ref{prop:xi-hellinger-toeplitz-symbol-extrema-qstar} 的 \(\sigma_\Delta(0)\) 与 \(\sigma_\Delta(\pi)\) 渐近给出。
有限主块的极端谱收敛为 Grenander--Szeg\H{o} 型结论；精度位数下界来自 \(\log_2\mathrm{cond}\) 的必要性。证毕。
\end{proof}

\begin{corollary}[精度--局域互补常数]\label{cor:xi-hellinger-toeplitz-precision-local-complementarity}
在命题 \ref{prop:xi-hellinger-toeplitz-conditioning-exponential-law} 的设定下，令
\[
a_{\mathrm{loc}}:=\frac{\Delta}{2}.
\]
则任意要求相对误差 \(\le\varepsilon\) 的稳定反演所需有效位数 \(B\) 满足
\[
B\ \gtrsim\ \frac{\pi^2}{a_{\mathrm{loc}}\ln 2}\ +\ O\bigl(1+\log(1/\varepsilon)\bigr).
\]
特别地，主项给出不可消去的互补下界
\[
\Bigl(\frac{B\ln 2}{\pi^2}\Bigr)\,a_{\mathrm{loc}}\ \gtrsim\ 1.
\]
\end{corollary}
\begin{proof}
将命题 \ref{prop:xi-hellinger-toeplitz-conditioning-exponential-law} 的主项下界
\(
B\gtrsim \frac{2\pi^2}{\Delta\ln 2}+O(1+\log(1/\varepsilon))
\)
代入 \(\Delta=2a_{\mathrm{loc}}\) 即得第一式；第二式为等价改写。证毕。
\end{proof}

\begin{theorem}[稠密极限的“面积律”自由能：\texorpdfstring{$f(\Delta)$}{f(Delta)} 的可审计模指数主项]\label{thm:xi-hellinger-toeplitz-free-energy-area-law}
\leanverified{paper\_xi\_hellinger\_toeplitz\_free\_energy\_area\_law}
设 \(\Delta>0\) 固定，并定义 Toeplitz 自由能密度
\[
f(\Delta):=\lim_{\kappa\to\infty}\frac{1}{\kappa}\log\det G_\kappa(\Delta)
=\frac{1}{2\pi}\int_{-\pi}^{\pi}\log\sigma_\Delta(\theta)\,\dd\theta
\]
（极限存在性见定理 \ref{thm:xi-hellinger-toeplitz-szego-thermo-limit}）。
则当 \(\Delta\downarrow0\) 时存在可审计的模指数展开
\[
\boxed{\;
f(\Delta)= -\frac{\pi^{2}}{\Delta}\ +\ \log\!\Bigl(\frac{\pi^{2}}{\Delta}\Bigr)\ +\ 2\log 2\ -\ \frac{\Delta}{16}\ +\ O\!\bigl(e^{-2\pi^{2}/\Delta}\bigr).\;}
\]
从而体积熵密度 \(-f(\Delta)\) 的主导项为 \(\pi^2/\Delta\)（面积律项），并伴随 \(\log(\pi^2/\Delta)\) 的边界对数修正与严格线性微校正 \(\Delta/16\)。
\end{theorem}
\begin{proof}[证明要点]
由命题 \ref{prop:xi-hellinger-toeplitz-symbol-poisson} 的 \(\vartheta_4\) 闭式与注 \ref{rem:xi-hellinger-toeplitz-modular-duality-scale} 的模对偶尺度，
\(\log\sigma_\Delta\) 的平均值可通过 Jacobi 乘积与 \(\tau\mapsto-1/\tau\) 变换转写为对偶 nome \(q_\star=e^{-2\pi^2/\Delta}\) 的级数。
对该级数在 \(q_\star\) 与 \(\Delta\) 的复合尺度下作显式展开，即得到所示主项与误差封口。证毕。
\end{proof}

\begin{proposition}[簇分解的指数可加性：深度体积熵在大间隙下呈分块加法]\label{prop:xi-hellinger-gram-cluster-additivity-exp}
将深度点集分为两簇 \(S^{(1)}\sqcup S^{(2)}\)，记 \(|S^{(\ell)}|=\kappa_\ell\)，并设簇间最小间隙
\[
L:=\min\{|x-y|:\ x\in S^{(1)},\ y\in S^{(2)}\}.
\]
对应 Gram 分块
\[
G=\begin{pmatrix}G_1 & C\\ C^{\mathsf T}& G_2\end{pmatrix},
\qquad
|C_{ij}|\le A(L)\le L\,e^{-L/2}.
\]
若 \(G_1,G_2\succ0\)，则存在常数 \(K\)（仅依赖 \(\|G_1^{-1}\|_2,\|G_2^{-1}\|_2\) 与 \(\kappa_1,\kappa_2\)）使
\[
\bigl|\log\det G-\log\det G_1-\log\det G_2\bigr|
\le K\,\|C\|_{\mathrm F}^{2}
\le K\,\kappa_1\kappa_2\,L^{2}e^{-L}.
\]
因此深度体积熵 \(S_{1/2}=-\log\det G\) 在簇间距 \(L\) 上具有严格的指数可加性：相互作用项的自然尺度为 \(e^{-L}\)。
\end{proposition}
\begin{proof}[证明要点]
由 Schur 补公式，
\[
\det G=\det G_1\cdot \det(G_2-C^{\mathsf T}G_1^{-1}C)
=\det G_1\det G_2\cdot \det(I-X),
\]
其中 \(X:=G_2^{-1/2}C^{\mathsf T}G_1^{-1}C\,G_2^{-1/2}\succeq 0\)。
当 \(\|X\|_2<1\) 时，
\(
|\log\det(I-X)|\le \mathrm{tr}(X)/(1-\|X\|_2)
\)。
并且
\[
\mathrm{tr}(X)=\|G_1^{-1/2}C\,G_2^{-1/2}\|_{\mathrm F}^2
\le \|G_1^{-1}\|_2\,\|G_2^{-1}\|_2\,\|C\|_{\mathrm F}^2,
\]
同时 \(\|X\|_2\le \|G_1^{-1}\|_2\|G_2^{-1}\|_2\|C\|_2^2\)。
将这些界并入常数 \(K\) 即得第一式。
最后一式由 \(\|C\|_{\mathrm F}^2\le \kappa_1\kappa_2\,\max_{i,j}|C_{ij}|^2\) 与 \(A(L)\le L e^{-L/2}\) 推出。证毕。
\end{proof}



\begin{proposition}[logistic 噪声下的多类别最优判决：最近邻规则与显式平均错误率]\label{prop:xi-logistic-multiclass-map-error}
令 \(\kappa\ge 2\)，并取有序深度对数谱
\(
\sigma_1<\sigma_2<\cdots<\sigma_\kappa
\)。
令离散输入 \(J\sim\mathrm{Unif}\{1,\dots,\kappa\}\)，观测信道为
\[
S=\sigma_J+U,
\qquad
U\sim\phi.
\]
则在等先验下最优 MAP（等价于最大似然）判决为最近邻规则
\[
\widehat J(S)\in\arg\min_{1\le j\le\kappa}|S-\sigma_j|,
\]
其决策边界为相邻中点 \((\sigma_j+\sigma_{j+1})/2\)。
并且等先验平均 Bayes 最小错误率具有严格闭式
\[
\boxed{\;
P_{\mathrm e}^{(\kappa)}
=\mathbb P(\widehat J(S)\neq J)
=\frac{2}{\kappa}\sum_{j=1}^{\kappa-1}\frac{1}{1+\exp\!\bigl((\sigma_{j+1}-\sigma_j)/2\bigr)}\; }.
\]
因此多类别可分性完全由相邻间隔 \(\Delta_j:=\sigma_{j+1}-\sigma_j\) 的 logistic 尾项决定。
\end{proposition}
\begin{proof}
因为 \(\phi\) 为偶函数且在 \([0,\infty)\) 上严格递减，故对固定 \(s\) 最大化 \(j\mapsto \phi(s-\sigma_j)\) 等价于最小化 \(j\mapsto |s-\sigma_j|\)，从而得到最近邻判决与中点边界。
\par
令 \(\Delta_j=\sigma_{j+1}-\sigma_j\)。
当 \(J=j\)（\(1<j<\kappa\)）时，误判事件等价于
\(
U\ge \Delta_j/2
\)
或
\(
U\le -\Delta_{j-1}/2
\)；
端点类 \(j=1,\kappa\) 仅有一侧尾部。
对称性给出
\(
\mathbb P(U\le -x)=\mathbb P(U\ge x)=1/(1+e^{x})
\)（\(x\ge 0\)），
对 \(j\) 求平均并将每个间隔的尾部贡献计数两次，即得所示闭式。证毕。
\end{proof}

% (User input window)

\begin{corollary}[由局部边界质量可逆重建深度间隔：错误剖面与间隔剖面一一对应]\label{cor:xi-logistic-boundary-mass-inverts-spacing}
\leanverified{paper\_xi\_logistic\_boundary\_mass\_inverts\_spacing}
在命题 \ref{prop:xi-logistic-multiclass-map-error} 的设定下，令相邻间隔 \(\Delta_j:=\sigma_{j+1}-\sigma_j>0\)。
定义局部边界质量（单侧尾概率）
\[
e_j:=\mathbb P\bigl(U\ge \Delta_j/2\bigr)=\frac{1}{1+e^{\Delta_j/2}}\in(0,1/2).
\]
则映射 \(\Delta_j\mapsto e_j\) 严格单调且显式可逆：
\[
\boxed{\;
\Delta_j=2\log\Bigl(\frac{1-e_j}{e_j}\Bigr).\;}
\]
因此一旦可观测到每个相邻切割的尾质量 \((e_j)_{j=1}^{\kappa-1}\)，深度间隔序列 \((\Delta_j)\) 可逐点无损恢复；并且
\(
P_{\mathrm e}^{(\kappa)}=\frac{2}{\kappa}\sum_{j=1}^{\kappa-1}e_j
\)
给出总错误率作为局部边界质量的线性汇总。
\end{corollary}
\begin{proof}
由 \(\mathbb P(U\ge x)=1/(1+e^{x})\)（\(x\ge 0\)）直接得到 \(e_j=(1+e^{\Delta_j/2})^{-1}\)，并由代数变形得到反演公式。
最后一式是命题 \ref{prop:xi-logistic-multiclass-map-error} 的误差闭式在记号 \(e_j\) 下的重写。证毕。
\end{proof}

\begin{corollary}[固定跨度下的最优深度布局：等距为误差极小解与必要跨度律]\label{cor:xi-logistic-optimal-spacing-span-law}
\leanverified{paper\_xi\_logistic\_optimal\_spacing\_span\_law}
在命题 \ref{prop:xi-logistic-multiclass-map-error} 的设定下，令总跨度
\[
L:=\sigma_\kappa-\sigma_1=\sum_{j=1}^{\kappa-1}\Delta_j.
\]
则有严格下界
\[
\boxed{\;
P_{\mathrm e}^{(\kappa)}
\ \ge\
\frac{2(\kappa-1)}{\kappa}\cdot\frac{1}{1+\exp\!\bigl(L/(2(\kappa-1))\bigr)}\; }.
\]
并且等号成立当且仅当
\(
\Delta_1=\cdots=\Delta_{\kappa-1}=L/(\kappa-1)
\)，即深度对数谱等距排列为唯一的误差极小解。
从而若要求 \(P_{\mathrm e}^{(\kappa)}\le \varepsilon\in(0,1)\)，则必有
\[
\boxed{\;
L\ \ge\ 2(\kappa-1)\log\!\Bigl(\frac{2(\kappa-1)}{\kappa\varepsilon}-1\Bigr)
=2(\kappa-1)\log\frac{1}{\varepsilon}+O(\kappa)\qquad(\varepsilon\downarrow0).\;}
\]
\end{corollary}
\begin{proof}
由命题 \ref{prop:xi-logistic-multiclass-map-error}，
\(
P_{\mathrm e}^{(\kappa)}=\frac{2}{\kappa}\sum_{j=1}^{\kappa-1}g(\Delta_j)
\)
其中 \(g(x)=(1+e^{x/2})^{-1}\)。
对 \(x>0\)，函数 \(g\) 可写为 \(g(x)=\bigl(1+e^{x/2}\bigr)^{-1}\) 并满足 \(g''(x)>0\)，故严格凸。
由 Jensen 不等式得到
\(
\frac{1}{\kappa-1}\sum_{j}g(\Delta_j)\ge g(L/(\kappa-1))
\)，
从而得到误差下界；严格性与等号条件给出唯一性。
跨度律由下界反解 \(L\) 得到。证毕。
\end{proof}

\begin{proposition}[互信息的熵差表示与低分辨率二次律：常数 \texorpdfstring{$1/6$}{1/6} 的刚性出现]\label{prop:xi-logistic-mutual-information-low-resolution}
在命题 \ref{prop:xi-logistic-multiclass-map-error} 的设定下，令输出混合密度
\[
p(s)=\frac{1}{\kappa}\sum_{j=1}^{\kappa}\phi(s-\sigma_j),
\]
并以自然对数定义微分熵 \(h(f):=-\int_{\RR}f\log f\)。
则互信息满足严格恒等式
\[
\boxed{\;I(J;S)=h(p)-h(\phi)=h(p)-2.\;}
\]
并且在低分辨率 regime（例如令中心化输入 \(X:=\sigma_J-\frac{1}{\kappa}\sum_{j}\sigma_j\) 且 \(\mathbb E[X^{2}]\) 足够小）存在二次近似律
\[
\boxed{\;
I(J;S)=\frac{1}{6}\,\Var(\sigma_J)+O\!\bigl(\mathbb E[|X|^{4}]\bigr).\;}
\]
其中系数 \(1/6\) 由核 \(\phi\) 的位置 Fisher 信息常数 \(\mathcal I_{\mathrm{loc}}(\phi)=1/3\) 强制确定，独立于 \(\kappa\) 与具体构型。
\end{proposition}
\begin{proof}
加性噪声信道给出 \(I(J;S)=h(S)-h(U)=h(p)-h(\phi)\)。
对 \(h(\phi)\)，作代换 \(x=e^{s}\) 得 \(\phi(s)\dd s=\dd x/(1+x)^{2}\) 且
\(
\log\phi(s)=\log x-2\log(1+x)
\)；
由 \(x\mapsto 1/x\) 的对称性可得
\(
\int_{0}^{\infty}\frac{\log x}{(1+x)^{2}}\dd x=0
\)，并且
\(
\int_{0}^{\infty}\frac{\log(1+x)}{(1+x)^{2}}\dd x=1
\)，故 \(h(\phi)=2\)。
\par
二次近似律属于加性噪声信道的小扰动展开：当输入尺度趋于零时，
\(
I(J;S)=\frac{\mathcal I_{\mathrm{loc}}(\phi)}{2}\Var(\sigma_J)+O(\mathbb E[|X|^4])
\)。
对 logistic 核，\(\phi'(u)/\phi(u)=-\tanh(u/2)\)，故
\(
\mathcal I_{\mathrm{loc}}(\phi)=\mathbb E[\tanh^{2}(U/2)]
\)。
令 \(T=\tanh(U/2)\in(-1,1)\)，则代换 \(u=2\,\operatorname{artanh}(t)\) 给出
\[
\phi(u)\,\dd u=\frac{1}{4\cosh^{2}(u/2)}\,\dd u=\frac12\,\dd t,
\]
从而 \(T\sim\mathrm{Unif}(-1,1)\) 且 \(\mathcal I_{\mathrm{loc}}(\phi)=\mathbb E[T^{2}]=1/3\)，得到系数 \(1/6\)。证毕。
\end{proof}

\leanverified{paper\_xi\_logistic\_fano\_lower\_bound\_from\_error}
\begin{corollary}[由显式平均错误率导出的 Fano 型可计算下界]\label{cor:xi-logistic-fano-lower-bound-from-error}
在命题 \ref{prop:xi-logistic-multiclass-map-error} 的设定下，有
\[
I(J;S)\ \ge\ \log\kappa\ -\ h_{\mathrm b}\!\bigl(P_{\mathrm e}^{(\kappa)}\bigr)\ -\ P_{\mathrm e}^{(\kappa)}\log(\kappa-1),
\]
其中 \(h_{\mathrm b}(u)=-u\log u-(1-u)\log(1-u)\) 为二元熵函数。
\end{corollary}
\begin{proof}
由 Fano 不等式。证毕。
\end{proof}

\begin{proposition}[位置 Fisher 信息的严格亏损：由后验方差给出的深度分散度证书]\label{prop:xi-logistic-fisher-information-deficit}
仍取
\(
p(s)=\frac{1}{\kappa}\sum_{j=1}^{\kappa}\phi(s-\sigma_j)
\)，并定义位置族 \(p_\theta(s)=p(s-\theta)\) 的 Fisher 信息
\[
\mathcal I_{\mathrm{loc}}(p):=\int_{\RR}\frac{(p'(s))^{2}}{p(s)}\,\dd s.
\]
令 \(J\sim\mathrm{Unif}\{1,\dots,\kappa\}\)、\(S=\sigma_J+U\)（\(U\sim\phi\)），并记
\[
T:=\tanh\!\Bigl(\frac{S-\sigma_J}{2}\Bigr),
\qquad
m(S):=\mathbb E[T\mid S].
\]
则存在严格恒等式
\[
\boxed{\;
\mathcal I_{\mathrm{loc}}(p)=\mathbb E[m(S)^{2}]
=\frac{1}{3}-\mathbb E\!\bigl[\Var(T\mid S)\bigr]
\ \le\ \frac{1}{3}\; }.
\]
并且等号成立当且仅当 \(\sigma_1=\cdots=\sigma_\kappa\)。
因此 \( \frac{1}{3}-\mathcal I_{\mathrm{loc}}(p)\) 是尺度曲率混合所引入的不可逆信息亏损，并以条件方差形式给出无拟合的可审计分散度证书。
\end{proposition}
\begin{proof}
由 \(\phi'(x)=-\tanh(x/2)\phi(x)\) 得
\[
p'(s)=-\frac{1}{\kappa}\sum_{j=1}^{\kappa}\tanh\!\Bigl(\frac{s-\sigma_j}{2}\Bigr)\phi(s-\sigma_j).
\]
因此
\(
-p'(s)/p(s)=\mathbb E[T\mid S=s]=m(s)
\)，从而
\(
\mathcal I_{\mathrm{loc}}(p)=\mathbb E[m(S)^{2}]
\)。
又因 \(T=\tanh(U/2)\)，同上代换 \(u=2\,\operatorname{artanh}(t)\) 得 \(T\sim\mathrm{Unif}(-1,1)\)，故 \(\mathbb E[T^{2}]=1/3\)（亦即 \(\mathcal I_{\mathrm{loc}}(\phi)=1/3\)），
由全方差分解
\(
\mathbb E[T^{2}]=\mathbb E[m(S)^{2}]+\mathbb E[\Var(T\mid S)]
\)
得到所示恒等式与上界。
若等号成立，则 \(\Var(T\mid S)=0\) 几乎处处，迫使 \(T\) 为 \(S\) 的确定函数；
在均匀先验的有限混合下，这等价于全部平移参数相同。证毕。
\end{proof}

\begin{remark}[信息几何的一维坍缩：以 \(\mathrm{TV}\) 为坐标的初等参数化]\label{rem:xi-logistic-tv-coordinate-parametrization}
令 \(v:=\|\phi_a-\phi_b\|_{\mathrm{TV}}\in[0,1)\)。
由命题 \ref{prop:xi-logistic-divergence-dictionary} 有严格反演
\[
|\Delta|=4\,\operatorname{artanh}(v)=2\log\!\Bigl(\frac{1+v}{1-v}\Bigr).
\]
并且
\[
H^{2}(\phi_a,\phi_b)
=1-\frac{(1-v^{2})\,\operatorname{artanh}(v)}{v},
\qquad
D_{\mathrm{KL}}(\phi_a\|\phi_b)
=2\frac{1+v^{2}}{v}\operatorname{artanh}(v)-2,
\]
\[
\chi^{2}(\phi_a\|\phi_b)
=\frac{16}{3}\frac{v^{2}}{(1-v^{2})^{2}},
\qquad
D_{2}(\phi_a\|\phi_b)
=\log\!\bigl(1+\chi^{2}(\phi_a\|\phi_b)\bigr).
\]
因此该平移族的若干核心散度可由单一标量坐标 \(v\) 完全参数化。
\end{remark}

\begin{remark}[Fisher--Rao 度量的显式化与二次近似偏离函数]\label{rem:xi-logistic-fisher-rao-and-curvature}
尺度核平移族 \(\{\phi_a\}\) 的位置 Fisher 信息为常数 \(\mathcal I_{\mathrm{loc}}(\phi)=1/3\)，故 Fisher--Rao 度量在参数 \(a\) 上为常系数，
任意两点的 Fisher--Rao 距离满足
\[
d_{\mathrm{FR}}(\phi_a,\phi_b)=\sqrt{\mathcal I_{\mathrm{loc}}(\phi)}\,|a-b|=\frac{|\Delta|}{\sqrt3}.
\]
并且由命题 \ref{prop:xi-logistic-divergence-dictionary} 的闭式可定义偏离函数
\[
\Psi(\Delta):=D_{\mathrm{KL}}(\phi_a\|\phi_b)-\frac{1}{2}d_{\mathrm{FR}}(\phi_a,\phi_b)^{2}
=|\Delta|\coth(|\Delta|/2)-2-\frac{\Delta^{2}}{6},
\]
其在 \(\Delta\to 0\) 时满足 \(\Psi(\Delta)=O(\Delta^{4})\)，并刻画 KL 与 Fisher--Rao 二次型在非线性可见尺度上的偏离。
\end{remark}

\begin{corollary}[多深度判别的边界质量表示：平均错误率的相邻切割分解]\label{cor:xi-logistic-boundary-mass-error}
在命题 \ref{prop:xi-logistic-multiclass-map-error} 的设定下，令 \(\Delta_j=\sigma_{j+1}-\sigma_j\)。
则平均错误率具有严格边界质量形式
\[
P_{\mathrm e}^{(\kappa)}
=\frac{2}{\kappa}\sum_{j=1}^{\kappa-1}\mathbb P\!\Bigl(U\ge \frac{\Delta_j}{2}\Bigr),
\qquad
\mathbb P(U\ge x)=\frac{1}{1+e^{x}}\quad(x\ge 0).
\]
因此 logistic 尺度核下的多类别误差不存在高维干涉项，仅由相邻间隔诱导的边界尾质量线性叠加给出。
\end{corollary}
\begin{proof}
由命题 \ref{prop:xi-logistic-multiclass-map-error} 的推导，误判等价于噪声落入相邻中点切割的两侧外尾，并由对称性化为单侧尾概率之和。证毕。
\end{proof}

\begin{remark}[\texorpdfstring{$4$}{4} 阶闭合阈值的散度表征]\label{rem:xi-logistic-fourth-order-closure}
命题 \ref{prop:xi-logistic-divergence-dictionary} 的小间隔展开显示：\(\mathrm{TV}\) 首项为线性（系数 \(1/4\)），而 \(H^{2}\) 与 \(D_{\mathrm{KL}}\) 的首个非平凡项为二次；
并且 \(D_{\mathrm{KL}}=4H^{2}+O(\Delta^{4})\) 把二者在四阶精度内锁定为同一局部几何。
因此当同时关注边界切割（由 \(\mathrm{TV}\) 的线性尾控制）与熵势涨落（由 \(H^{2}\)、\(D_{\mathrm{KL}}\) 的二次项控制）时，四阶项成为最小的闭合阶。
\end{remark}

\begin{corollary}[以信息可见性为目标的深度谱设计律：最优构型收缩到等距族]\label{cor:xi-logistic-depth-design-law}
在固定 \(\kappa\) 与总跨度 \(L\) 的约束下，使命题 \ref{prop:xi-logistic-multiclass-map-error} 的平均错误率 \(P_{\mathrm e}^{(\kappa)}\) 最小的构型唯一地由等距谱实现：
\[
\sigma_j=\sigma_1+\frac{j-1}{\kappa-1}L\qquad(1\le j\le\kappa).
\]
并且该构型同时使推论 \ref{cor:xi-logistic-fano-lower-bound-from-error} 给出的互信息下界达到最大值。
\end{corollary}
\begin{proof}
由推论 \ref{cor:xi-logistic-optimal-spacing-span-law}，在固定跨度下 \(P_{\mathrm e}^{(\kappa)}\) 的下界由 Jensen 取等条件唯一实现；
而 Fano 下界对 \(P_{\mathrm e}^{(\kappa)}\) 单调递减，故同一构型也最大化该下界。证毕。
\end{proof}
